\documentclass[11pt,a4paper]{article}


\usepackage[utf8]{inputenc}
\usepackage{amsmath,amssymb}
\usepackage{mathtools}
\usepackage{graphicx}
\usepackage{booktabs}
\usepackage[margin=1in]{geometry}
\usepackage{authblk}
\usepackage{siunitx}
\usepackage[hidelinks]{hyperref}
\usepackage{caption}
\usepackage{xcolor}
\usepackage{microtype}
\usepackage{csquotes}
\usepackage[numbers,sort&compress]{natbib}
% ---- float placement: these figures are tall (stacked two-panel portrait),
% so the default \topfraction of 0.7 would defer them to the end of the document.
\usepackage[section]{placeins}
\renewcommand{\topfraction}{0.92}
\renewcommand{\bottomfraction}{0.70}
\renewcommand{\textfraction}{0.05}
\renewcommand{\floatpagefraction}{0.72}
\setcounter{topnumber}{3}
\setcounter{bottomnumber}{2}
\setcounter{totalnumber}{5}
% Floats are confined to their own section (placeins), but NOT to their own
% subsection: barriering at subsection level forced page breaks that left
% large gaps. A table may therefore appear just above the subsection heading
% that discusses it; that is standard float behaviour and is preferred here to
% half-empty pages.



\allowdisplaybreaks
\setlength{\parindent}{15pt}
\setlength{\parskip}{0.3em}
\linespread{1.03}

\captionsetup{font=small,labelfont=bf}

\graphicspath{{.}{./figures/}{./fig/}{./images/}}

% ---- shared macros (consistent with Papers I--III) ----
\newcommand{\gmean}{\langle g\rangle}
\newcommand{\RA}{R_{A}}
\newcommand{\RV}{R_{V}}
\newcommand{\Ptwo}{P_{2}}
\newcommand{\Pthree}{P_{3}}

% ---- Paper IV macros: spatial gravity, slope, geopotential dispersion ----
\newcommand{\geff}{\mathbf{g}_{\mathrm{eff}}}
\newcommand{\ngrav}{\mathbf{g}}
\newcommand{\acf}{\mathbf{a}_{\mathrm{c}}}
\newcommand{\nhat}{\hat{\mathbf{n}}}
\newcommand{\slope}{\theta}
\newcommand{\slopeavg}{\bar{\slope}}
\newcommand{\Upot}{U}
\newcommand{\Ubar}{\bar{U}}
\newcommand{\potvar}{\sigma_{U}}
\newcommand{\potvarn}{\tilde{\sigma}_{U}}
\newcommand{\omegarot}{\omega}
\newcommand{\dens}{\varrho}
\newcommand{\Nf}{N_{f}}



\begin{document}
\emergencystretch=2em
% ==================== TITLE BLOCK ====================
\begin{center}
{\LARGE\bfseries Spatial surface gravity of small bodies from polyhedral
shape models: slope, geopotential, and the resolution robustness of
geopotential dispersion\par}
\vspace{0.8em}
{\itshape Prapon Kanjanatarayont\par}
\vspace{0.2em}
{\small\color{gray} Independent Researcher \textperiodcentered\ Thailand\par}
\vspace{0.1em}
{\small\color{gray} \texttt{prapon.kanjana@gmail.com}\par}
\end{center}
\vspace{1.0em}




{\noindent\large\bfseries Abstract}\\[0.4em]
{\noindent\small
Three companion papers established that the area-weighted mean surface gravity
$\gmean$ of a small body can be estimated from a polyhedral shape model,
without an independently determined density (Papers~I--III). A single scalar,
however, cannot say \emph{where} on a body the gravity is high or low, how
steeply the surface is inclined to the local effective gravity, or where loose
regolith is driven to accumulate. Here we extend the shape-based framework from
the scalar mean to the spatially resolved surface fields: the effective gravity
vector $\geff$ (self-gravitation plus rotation), the gravitational slope
$\slope$, and the surface geopotential $\Upot$, evaluated facet by facet for
seven asteroids from a polyhedral shape model together with an assumed uniform
density and a specified spin state.

The framework is validated against an independent per-facet reference. For
(162173)~Ryugu we compare with the gravitational products released by the
Hayabusa2 team on the identical $49{,}152$-facet shape model at the same
density and rotation period. Agreement is essentially exact: mean slope
$15.06^{\circ}$ against $15.06^{\circ}$, RMS difference $0.48^{\circ}$,
$99.96\%$ of facets within $5^{\circ}$ and $99.95\%$ within $2^{\circ}$,
Pearson $r=0.9989$ in slope and $0.9898$ in geopotential. For (101955)~Bennu the global
mean slope is $15.5^{\circ}$ and the latitudinal profile rises from
$10.5^{\circ}$ near the equator to $19.9^{\circ}$ at mid-latitude before
falling toward the pole, reproducing the pattern reported from OSIRIS-REx data.

The principal methodological result concerns stability under mesh refinement.
Because slope is a \emph{local} quantity, defined through individual facet
normals, the area-weighted mean slope varies by $1.8\%$ to $25.6\%$ across
meshes of the same body decimated from roughly $42{,}000$ to $2{,}000$ facets,
whereas the normalized geopotential dispersion $\potvarn$ varies by only $0.9\%$
to $5.9\%$, and is the more stable of the two on every body tested. The
sharper distinction is stabilization: between the two finest meshes the
geopotential dispersion is already near a plateau (changes of $0.00\%$ to
$0.60\%$), whereas the mean slope shows no comparable plateau within the tested
resolution range, still rising by up to $8.2\%$. A slope value is therefore not
reproducible unless the mesh resolution is quoted with it, whereas the
geopotential dispersion exhibits practical convergence over the tested range.

Comparing two independent shape models of Itokawa---the spacecraft-derived
Gaskell model and a ground-based radar model---separates what survives a coarser
shape from what does not. Bulk descriptors are compared at a common facet count
by interpolation, and distribution statistics at the nearest available direct
meshes. Bulk descriptors agree
closely (area-weighted mean slope within $0.8\%$, geopotential dispersion within
$7.9\%$), but the distribution tails do not: the radar model reaches only
$41^{\circ}$ with $1.7\%$ of facets above $30^{\circ}$, against $149^{\circ}$
and $4.8\%$ for the Gaskell model. In this controlled case a ground-based shape
therefore preserves the global gravity descriptors while failing to reproduce
the distribution tails; one case does not establish the result generally, but it
indicates where coarse models may and may not be relied upon. We verify the computed fields against two exact solutions: the
closed-form potential of a homogeneous ellipsoid (Pearson $r=1.000000$, mean
relative error $0.18\%$) and the analytic slope field of a rotating sphere.
The latter test is essential---a non-rotating sphere check cannot detect an
error in the centrifugal term---and we recommend it as a standard verification
for surface-gravity codes. Both tests are released as runnable scripts together
with the analysis code.
}

\medskip
\noindent\textbf{Keywords:} Asteroids; Surface gravity; Gravitational slope;
Geopotential; Shape models; Polyhedral gravity

\section{Introduction}
\label{sec:intro}

The surface gravity of a small body governs nearly every process that shapes
its regolith and constrains how a spacecraft can operate on or near it. Where
material comes to rest, how loose regolith migrates, whether a boulder is
stable on a slope, and where a lander or sampling head can safely touch down
are all set by the local effective gravity---the vector sum of self-gravitation
and the centrifugal acceleration of rotation
\citep{scheeres2016,scheeres2019}. For the bodies visited by spacecraft these
quantities are now mapped in detail from mission data. For the far larger
population resolved only as a shape model---from ground-based radar, lightcurve
inversion, or a single flyby---no such maps exist, even though the same physics
applies.

The companion papers of this series developed a purely shape-based route to the
\emph{mean} surface gravity. Paper~I \citep{paper1} showed that the
area-equivalent radius $\RA$ outperforms the volume-equivalent radius $\RV$ for
elongated bodies and proved the bound $GM/\RA^{2}\leq\gmean$; Paper~II
\citep{paper2} derived the optimal blended radius from a non-analytic
$\Pthree/\Ptwo$ deformation law; and Paper~III \citep{paper3} tested that
estimator against a homogeneous-polyhedron reference field for $22$ non-convex
models, obtaining a mean absolute error of $1.72\%$, and delimited its domain of
validity. That work establishes that one scalar---the area-weighted mean
gravity---can be recovered from the shape model to within a few per cent. A
scalar
mean, however, is silent about the spatial structure of the field. Those
questions require moving from a single number to a field defined over the
entire surface.

This paper makes that step. Using the same polyhedral shape models and an
assumed uniform density, we compute three spatially resolved quantities at
every facet: the effective gravity $\geff$, the gravitational slope $\slope$
(the angle between $\geff$ and the outward facet normal $\nhat$), and the
geopotential $\Upot$, whose minima mark the sinks toward which loose material
migrates. The self-gravitation is evaluated with the analytic polyhedral method
of \citet{werner1996}, as implemented in a modern open-source library
\citep{tsoulis2012,gravitylib}, so that the framework remains light enough to
run on a laptop yet remains accurate on the concave, bilobed, and highly
irregular shapes that dominate the small-body population.

Making these fields easy to compute from a shape model, an assumed density and
a specified spin state is useful, but it also
exposes a methodological hazard that, to our knowledge, has not been stated
plainly in the shape-based context: the slope field is sensitive to the
resolution of the shape model. Because slope is a local quantity, defined
through the orientation of each individual facet, refining a mesh changes the
distribution of facet normals and hence the mean slope, even when the underlying
body is identical. Section~\ref{sec:slopesens} shows that over a
$2{,}000$-to-$42{,}000$-facet range the area-weighted mean slope varies by up to
a quarter of its own value, and---more tellingly---that it has still not
converged at the finest meshes available. A slope value quoted without its mesh
resolution is therefore not reproducible, a point with direct consequences for
how such numbers are reported and compared between studies.

The central positive result of this paper is that a closely related
quantity---the area-weighted dispersion of the surface geopotential---behaves far
better. Because the geopotential is an integral of the mass distribution rather
than a local surface property, the facet-scale roughness that scatters the slope
largely cancels in the potential. Across the bodies tested, $\potvarn$ is the
more stable of the two diagnostics in every case, and it approaches a practical
plateau over the tested refinement range whereas the slope shows no comparable
stabilization (Section~\ref{sec:varrobust}). This is
consistent with the behaviour noted by \citet{kanamaru2019} for Itokawa and
makes $\potvarn$ the natural descriptor to quote when a resolution-independent
statement about a body's surface-gravity field is required.

A separate question is what survives when the shape model itself is coarse, as
it must be for objects never visited by a spacecraft. Itokawa is the rare case
with two independent shape models of comparable extent---the spacecraft-derived
Gaskell model and a ground-based radar model---and comparing them at matched
resolution (Section~\ref{sec:crossversion}) separates the descriptors that
tolerate a ground-based shape from those that do not.

We validate the framework against bodies for which detailed gravitational
products already exist. For (162173)~Ryugu we compare directly against the
per-facet gravitational products released with \citet{kanamaru2019}, computed on
the identical shape model at the same density and rotation period; this is the
strongest available test, since shape, method, density, and spin are all held
fixed. For (101955)~Bennu the recovered mean slope and latitudinal profile are
compared with the OSIRIS-REx results
\citep{scheeres2019,barnouin2020,jawin2020}. We then use (433)~Eros,
(25143)~Itokawa, (951)~Gaspra, (243)~Ida, and (216)~Kleopatra to characterise
the resolution behaviour of the two diagnostics across the range of small-body
morphologies \citep{zuber2000,thomas2002,fujiwara2006,abe2006}.

\medskip
Three methodological commitments govern the analysis.

First, \emph{every slope figure is reported with its mesh resolution}. Because
the sensitivity documented here is a property of the diagnostic rather than of
the body, a slope value quoted without a facet count is not reproducible. All
slope numbers in this paper are accompanied by the facet count at which they
were obtained. Cross-model comparisons are made at matched resolution wherever
the meshes permit it, and where they do not---as for the tails of the Itokawa
distributions---the facet counts actually used are stated.

Second, \emph{the fields are verified against exact solutions, including one
that exercises the rotating frame}. The polyhedral potential is compared with
the closed-form potential of a homogeneous ellipsoid, and the full slope field
with the analytic slope of a uniformly rotating sphere. The second test is not
redundant: a non-rotating sphere reproduces zero slope for either sign of the
centrifugal term, so it cannot detect an error there, whereas the rotating case
discriminates sharply. We report both, and recommend the rotating test as a
standard check for codes of this kind.

Third, \emph{the uniform-density assumption is stated rather than hidden}. All
fields in this paper are computed for a homogeneous interior. This assumption
holds exactly for none of the bodies considered, and the resulting
limitations---in particular for Itokawa, whose head and body are known to
differ in density---are examined in Section~\ref{sec:discussion} rather than
deferred.

The paper is organised as follows. Section~\ref{sec:methods} defines the
effective gravity, slope, and geopotential, describes the polyhedral solver and
the shape models used, and specifies the decimation and smoothing operators.
Section~\ref{sec:validation} presents the validation against Ryugu and Bennu.
Section~\ref{sec:slopesens} documents the resolution behaviour of the slope,
and Section~\ref{sec:varrobust} that of the geopotential dispersion.
Section~\ref{sec:crossversion} compares independent shape versions of the same
body. Section~\ref{sec:discussion} discusses the consequences for cross-study
comparison and for interior-structure inference, together with the limitations
of the uniform-density assumption, and Section~\ref{sec:conclusions} concludes.

\section{Data and methods}
\label{sec:methods}

This section defines the three surface fields computed throughout the paper,
fixes the sign conventions that make them reproducible, describes the
polyhedral solver and the shape models, specifies the mesh operators used to
probe resolution behaviour, and reports the verification of the implementation
against exact analytic solutions.

\subsection{Target quantities}
\label{sec:targets}

Let the body occupy a region $B$ bounded by a closed triangulated surface with
$\Nf$ facets, and let the interior density $\dens$ be uniform. The
gravitational potential at a field point $\mathbf{r}$ is

\begin{equation}
U_{g}(\mathbf{r}) \;=\; -\,G\dens \int_{B}
\frac{\mathrm{d}^{3}r'}{\lvert \mathbf{r}-\mathbf{r}' \rvert},
\label{eq:Ug}
\end{equation}

so that $U_{g}<0$ and the self-gravitational attraction is

\begin{equation}
\ngrav(\mathbf{r}) \;=\; -\nabla U_{g}(\mathbf{r}),
\label{eq:g}
\end{equation}

which points into the body. For a body rotating uniformly at angular rate
$\omegarot$ about the $z$ axis of the body-fixed frame, write
$\mathbf{r}_{\perp}=(x,y,0)$ for the component perpendicular to the spin axis.
The centrifugal acceleration is directed away from the spin axis,

\begin{equation}
\acf(\mathbf{r}) \;=\; \omegarot^{2}\,\mathbf{r}_{\perp},
\label{eq:acf}
\end{equation}

and the effective acceleration experienced by a particle at rest in the
rotating frame is

\begin{equation}
\geff(\mathbf{r}) \;=\; \ngrav(\mathbf{r}) + \acf(\mathbf{r}).
\label{eq:geff}
\end{equation}

The \emph{effective gravitational slope} $\slope$---hereafter simply the
gravitational slope, the qualifier being understood---of facet $i$ with outward
unit normal $\nhat_{i}$ is the angle between the upward direction $-\geff$ and
that normal,

\begin{equation}
\slope_{i} \;=\; \arccos\!\left(
\frac{-\,\geff(\mathbf{c}_{i})\cdot\nhat_{i}}
     {\lVert \geff(\mathbf{c}_{i}) \rVert} \right),
\label{eq:slope}
\end{equation}

evaluated at the facet centroid $\mathbf{c}_{i}$. A surface perpendicular to
the local effective gravity has $\slope=0$; a surface on which the effective
gravity has an outward component has $\slope>90^{\circ}$, which is possible
only where the centrifugal term dominates or where the local topography
overhangs.

The \emph{geopotential} adds the rotational term to the gravitational
potential,

\begin{equation}
\Upot(\mathbf{r}) \;=\; U_{g}(\mathbf{r})
\;-\;\tfrac{1}{2}\,\omegarot^{2}\,\lVert\mathbf{r}_{\perp}\rVert^{2},
\label{eq:geopot}
\end{equation}

so that $\geff=-\nabla\Upot$ and the minima of $\Upot$ on the surface mark the
sinks toward which loose material is driven.

Two scalar summaries of these fields are used throughout. Weighting each facet
by its area $A_{i}$, the area-weighted mean slope is

\begin{equation}
\slopeavg \;=\;
\frac{\sum_{i} A_{i}\,\slope_{i}}{\sum_{i} A_{i}},
\label{eq:slopeavg}
\end{equation}

and the area-weighted geopotential mean and variance are

\begin{equation}
\Ubar \;=\; \frac{\sum_{i} A_{i}\,\Upot_{i}}{\sum_{i} A_{i}},
\qquad
\potvar^{2} \;=\;
\frac{\sum_{i} A_{i}\,\bigl(\Upot_{i}-\Ubar\bigr)^{2}}{\sum_{i} A_{i}}.
\label{eq:potvar}
\end{equation}

The square root $\potvar$ of Eq.~\eqref{eq:potvar} is the corresponding
area-weighted standard deviation. Because $\Upot$ carries units and its zero
point is fixed by Eq.~\eqref{eq:Ug}, comparisons across bodies use the
dimensionless \emph{normalized geopotential dispersion}

\begin{equation}
\potvarn \;=\; \frac{\potvar}{\lvert \Ubar \rvert}.
\label{eq:potvarn}
\end{equation}

which is the coefficient of variation of the surface geopotential. We use the
term \emph{dispersion} rather than \emph{variance} throughout, because
$\potvarn$ is built from the standard deviation and not from
$\potvar^{2}$.

Both $\slopeavg$ and $\potvarn$ are area-weighted rather than facet-weighted,
so that neither is biased by the uneven facet areas that decimation produces.

\subsubsection{Sign conventions}
\label{sec:signs}

Equations~\eqref{eq:g}--\eqref{eq:geopot} fix every sign, but two conventions
in common use differ from them and are worth stating explicitly, because a
discrepancy in either is easy to introduce and hard to detect.

First, some solvers---including the library used here---return the positive
Newtonian integral $-U_{g}$ rather than $U_{g}$ itself, together with an
acceleration vector that points into the body. The implementation negates the
returned potential before forming Eq.~\eqref{eq:geopot}, so that the
gravitational and rotational terms of the geopotential carry consistent signs.

Second, the centrifugal term must \emph{oppose} self-gravity, as in
Eq.~\eqref{eq:acf}: at the equator of a rotating body the effective
acceleration is smaller in magnitude than the attraction alone. The opposite
sign is easy to adopt by mistake when the underlying solver stores gravity as
an outward-pointing vector, and---as Section~\ref{sec:numval} shows---it cannot
be detected by testing a non-rotating body.

\subsubsection{Surface-limit convention}
\label{sec:surflimit}

The Newtonian potential and its gradient possess well-defined surface limits
for a homogeneous solid of finite density carrying no surface mass layer: the
physical acceleration is finite and, away from edges and vertices, continuous
across a face. In a polyhedral implementation, however, evaluating exactly on a
face makes the solid-angle and edge terms of the analytic expression ambiguous
at the level of a branch choice. Fields are therefore evaluated at

\begin{equation}
\mathbf{p}_{i} \;=\; \mathbf{c}_{i} \;-\; \epsilon\,\nhat_{i},
\qquad
\epsilon \;=\; 10^{-4}\,L,
\label{eq:eps}
\end{equation}

where $L$ is the diagonal of the shape model's bounding box, so that the
evaluation point lies just inside the surface and the interior limit of the
numerical expression is selected consistently. The displacement is four orders
of magnitude below the body scale and two to three orders below a typical facet
dimension, and the results are insensitive to it: halving or doubling
$\epsilon$ changes $\slopeavg$ and $\potvarn$ in the fifth significant figure.

\subsection{Polyhedral field evaluation}
\label{sec:solver}

The self-gravitation of a homogeneous polyhedron is computed with the analytic
line-and-face-integral method of \citet{werner1996}, in the formulation of
\citet{tsoulis2012}, as implemented in the open-source \texttt{polyhedral-%
gravity} package \citep{gravitylib}. The method is exact for a uniform
polyhedron---no series truncation, no mascon approximation---and is valid for
non-convex and multiply concave surfaces, which is essential for the bilobed and
irregular bodies considered here. Facet normals are oriented outward before
evaluation, and the solver is run in its outward-normal mode.

Since Eqs.~\eqref{eq:Ug}--\eqref{eq:g} are linear in $\dens$, the field for any
assumed density follows from a single unit-density evaluation by scaling. This
is used to sweep the assumed bulk density at fixed shape
(Section~\ref{sec:validation}) at the cost of one solver call, and it motivates the
region-superposition strategy discussed as a future interior-density extension
in Section~\ref{sec:interior}.

All lengths in the shape files are in kilometres; internally they are converted
to SI, and slopes are reported in degrees.

\subsection{Shape models}
\label{sec:shapes}

Table~\ref{tab:shapes} lists the nine shape models used, covering seven
asteroids and spanning oblate top shapes, elongated and bilobed configurations,
angular and irregular fragments, and one extreme dog-bone. Two bodies are
represented by two independent models each, which permits the cross-version
comparison of Section~\ref{sec:crossversion}.

Bulk densities are the measured values where a spacecraft or a satellite has
constrained the mass, and representative taxonomic values otherwise; because
the fields scale with $\dens$ in a known way, an assumed density affects the
absolute slope but not the resolution behaviour that is the subject of
Sections~\ref{sec:slopesens} and~\ref{sec:varrobust}. For Ryugu the density is
set to $1200\ \mathrm{kg\,m^{-3}}$ exactly, matching the value used for the
reference products against which we validate, rather than to the marginally
different measured bulk value.

\begin{table}[t]
\centering
\small
\caption{Shape models used in this work. $\Nf$ is the facet count of the
delivered model. Densities are in \si{\gram\per\cubic\centi\metre} and
rotation periods in hours.}
\label{tab:shapes}
\begin{tabular}{llrrrl}
\toprule
Body & Model & $\Nf$ & $\dens$ & $P$ & Role \\
\midrule
(101955) Bennu     & OLA v20            & 196\,608 & 1.194 & 4.2961  & validation \\
(162173) Ryugu     & SFM 49k v20180804  &  49\,152 & 1.200 & 7.627   & validation \\
(433) Eros         & Gaskell            &  49\,152 & 2.670 & 5.270   & resolution \\
(25143) Itokawa    & Gaskell            &  49\,152 & 1.950 & 12.132  & resolution, cross-version \\
(25143) Itokawa    & radar-based        &   3\,688 & 1.950 & 12.132  & cross-version \\
(951) Gaspra       & Thomas             &  32\,040 & 2.700 & 7.042   & resolution \\
(243) Ida          & Thomas             &  32\,040 & 2.600 & 4.634   & resolution \\
(216) Kleopatra    & Ostro radar        &   4\,092 & 3.380 & 5.385   & cross-version \\
(216) Kleopatra    & \texttt{kleo.v2.7} &   2\,292 & 3.380 & 5.385   & cross-version \\
\bottomrule
\end{tabular}
\end{table}

Every model is checked before use. The mesh is verified to be watertight and
manifold, with no degenerate facets and no inward-pointing normals; the
principal axes of inertia are computed and the spin axis confirmed to coincide
with the axis of maximum inertia, so that the latitude used in the profiles of
Section~\ref{sec:validation} is meaningful; and the bounding box is compared
with the published dimensions of the body, which is the check that catches unit
errors and mis-scaled models. This last test is not a formality: it is what
identifies the two Kleopatra models as differing in scale as well as in detail
(Section~\ref{sec:crossversion}).

\subsection{Mesh operators}
\label{sec:operators}

Two operators are applied to the meshes, and they are deliberately different in
character.

\paragraph{Decimation.} To vary resolution while preserving the body, vertices
are clustered on a uniform grid: the bounding box is divided into cells, all
vertices falling in a cell are replaced by their centroid, faces that become
degenerate are removed, and duplicate faces are collapsed. The grid is refined
until the target vertex count is reached. This reduces $\Nf$ while leaving the
overall figure of the body essentially unchanged.

It is important that decimation does \emph{not} smooth the surface. A coarser
facet spans the same topography that several finer facets spanned before, so
its normal still reflects the local relief; what changes is how that relief is
sampled, not whether it is present. This distinction is what separates the
resolution effect of Section~\ref{sec:slopesens} from the roughness effect
below.

\paragraph{Smoothing.} To remove facet-scale relief while holding $\Nf$ fixed,
Laplacian smoothing is applied: each vertex is displaced toward the centroid of
its topological neighbours,

\begin{equation}
\mathbf{v}_{j} \;\leftarrow\; \mathbf{v}_{j}
+ \lambda\left( \frac{1}{\lvert N(j)\rvert}\sum_{m\in N(j)}\mathbf{v}_{m}
- \mathbf{v}_{j}\right),
\qquad \lambda = 0.5,
\label{eq:laplace}
\end{equation}

iterated a prescribed number of times. Successive iterations progressively
suppress short-wavelength relief while leaving the global figure of the body
intact. Applying the two operators separately makes it possible to attribute a
change in $\slopeavg$ either to how the surface is sampled (decimation) or to
how rough it is (smoothing).

\subsection{Resolution diagnostics}
\label{sec:diagnostics}

For a scalar diagnostic $X$ evaluated on a set of meshes with facet counts
$\Nf^{(1)}<\dots<\Nf^{(K)}$ obtained by decimating a single model, two measures
are reported.

The \emph{range sensitivity} is the peak-to-peak variation normalized by the
mean over the set,

\begin{equation}
S[X] \;=\; \frac{\max_{k} X^{(k)} - \min_{k} X^{(k)}}
{\tfrac{1}{K}\sum_{k} X^{(k)}}.
\label{eq:sens}
\end{equation}

The \emph{convergence residual} is the fractional change between the two finest
meshes,

\begin{equation}
C[X] \;=\;
\frac{\bigl\lvert X^{(K)} - X^{(K-1)} \bigr\rvert}{X^{(K-1)}}.
\label{eq:conv}
\end{equation}

Both are quoted as percentages. By construction $C[X]\ge 0$: it records the
\emph{magnitude} of the fractional change over the final refinement step. Where
the direction of that change is also of interest, the corresponding signed
quantity

\begin{equation}
C_{\mathrm{signed}}[X] \;=\;
\frac{X^{(K)} - X^{(K-1)}}{X^{(K-1)}}
\label{eq:convsigned}
\end{equation}

may be reported alongside it. Every sequence studied here increases
monotonically over its final refinement step, so the two coincide in magnitude
and $C_{\mathrm{signed}}$ is positive throughout.

The two answer different questions. $S[X]$ measures how much a reported value
would change had a coarser or finer mesh been used anywhere in the tested range,
and is the quantity relevant when comparing published numbers. $C[X]$ measures
whether the diagnostic is approaching a limit as the mesh is refined, and is
the quantity relevant to asking whether a resolution-independent value exists
at all. A diagnostic may have a modest $S$ merely because it varies slowly, yet
a large $C$ because it has not settled; the two are reported together
throughout.

\subsection{Numerical verification}
\label{sec:numval}

The implementation is verified against two closed-form solutions, chosen so
that between them they exercise every term in
Eqs.~\eqref{eq:Ug}--\eqref{eq:geopot}.

\paragraph{Potential: homogeneous ellipsoid.} For a uniform ellipsoid with
semi-axes $a\ge b\ge c$ the exterior potential is given exactly by the
classical elliptic integral

\begin{equation}
U_{g}(\mathbf{r}) \;=\; -\pi G\dens\,abc \int_{\lambda}^{\infty}
\left(1 - \frac{x^{2}}{a^{2}+u} - \frac{y^{2}}{b^{2}+u}
        - \frac{z^{2}}{c^{2}+u}\right)
\frac{\mathrm{d}u}{\Delta(u)},
\label{eq:ellipsoid}
\end{equation}

with $\Delta(u)=\sqrt{(a^{2}+u)(b^{2}+u)(c^{2}+u)}$ and $\lambda$ the largest
root of the ellipsoidal coordinate equation, which vanishes on the surface. In
the limit $a=b=c=R$ this reduces to $-GM/R$ on the surface, and the quadrature
used reproduces that limit exactly. Evaluated on a triangulated ellipsoid with
$a,b,c=1500,1000,800$~m and $\dens=2000\ \mathrm{kg\,m^{-3}}$, the polyhedral
potential agrees with Eq.~\eqref{eq:ellipsoid} facet by facet to a Pearson
correlation of $1.000000$, with a mean relative error of $0.18\%$ and a maximum
of $0.21\%$; the residual is consistent with the polygonal approximation of the
curved surface.

\paragraph{Slope: uniformly rotating sphere.} For a homogeneous sphere of
radius $R$ and surface gravity $g=GM/R^{2}$ rotating at rate $\omegarot$, write
$k=\omegarot^{2}R/g$ for the ratio of centrifugal to gravitational acceleration
at the equator. Combining Eqs.~\eqref{eq:acf}--\eqref{eq:slope} gives the slope
at latitude $\phi$ in closed form,

\begin{equation}
\cos\slope(\phi) \;=\;
\frac{1-k\cos^{2}\phi}
     {\sqrt{(1-k)^{2}\cos^{2}\phi+\sin^{2}\phi}}.
\label{eq:rotsphere}
\end{equation}

At $k=0.5$ this predicts $19.11^{\circ}$, $18.43^{\circ}$ and $13.90^{\circ}$
at latitudes $30^{\circ}$, $45^{\circ}$ and $60^{\circ}$; the implementation
returns $18.97^{\circ}$, $18.33^{\circ}$ and $13.91^{\circ}$, agreeing to
better than $0.15^{\circ}$ at the mesh resolution used.

This second test is not redundant, and the reason is worth stating. Reversing
the sign of $\acf$ in Eq.~\eqref{eq:geff} is equivalent to replacing $k$ by
$-k$ in Eq.~\eqref{eq:rotsphere}, which at $k=0.5$ changes the predicted slope
at $45^{\circ}$ latitude from $18.43^{\circ}$ to $11.31^{\circ}$---a
discrepancy of nearly two-fifths that the rotating test detects immediately.
A \emph{non-rotating} sphere, by contrast, returns zero slope for either sign,
and so cannot detect the error at all. Since a static sphere is the natural
first test to write, we recommend that codes of this kind be checked against
Eq.~\eqref{eq:rotsphere} as well; both tests are provided with the code.

As a control, the non-rotating sphere is also run and returns a mean slope of
$0.35^{\circ}$, consistent with the polygonal discretisation of a curved
surface at the resolution used.

\subsection{Reproducibility}
\label{sec:repro}

All computations use Python~3.13 with \texttt{numpy} and \texttt{scipy}, and
version~3.3.1 of the \texttt{polyhedral-gravity} package \citep{gravitylib}
for the self-gravitation. No random numbers enter the analysis: every quantity
reported is deterministic given the shape model, the assumed density, and the
rotation period.

The analysis code, the verification scripts implementing
Eqs.~\eqref{eq:ellipsoid} and~\eqref{eq:rotsphere}, the mesh diagnostics, and a
machine-readable file containing every numerical result quoted in this paper are
archived and openly available; sources for all shape models are listed there.
Shape models themselves are not redistributed.

\section{Validation}
\label{sec:validation}

The framework is tested against two bodies for which independent gravitational
products exist. Ryugu provides the stronger test, because the Hayabusa2 team
has released per-facet gravitational quantities computed on a shape model that
we can use unchanged, at a stated density and rotation period; the comparison
is therefore facet by facet with every input held fixed. Bennu provides a
complementary test at the level of published summary statistics, and adds a
check that the framework responds correctly to the assumed bulk density.

\subsection{Ryugu: comparison facet by facet}
\label{sec:ryugu}

\citet{kanamaru2019} and the archive accompanying the Hayabusa2 shape analysis
distribute, alongside the $49\,152$-facet SFM shape model of Ryugu, a table of
gravitational potential and gravitational slope for every plate of that model,
computed with a Werner--Scheeres polyhedral solver for a uniform density of
$1200\ \mathrm{kg\,m^{-3}}$ and a rotation period of $7.63$~h. Because shape,
method, density, and spin are all fixed by that specification, any disagreement
must originate in the implementation rather than in the inputs. This is the
strongest test available to us, and the only one in this paper that reaches
below the level of aggregate statistics.

It is worth being precise about what such a test does and does not establish.
The reference products are themselves computed for a homogeneous interior with a
Werner--Scheeres solver, so the comparison is an \emph{implementation}
validation: it confirms that our sign conventions, facet correspondence and
numerical evaluation agree with an independent implementation of the same
physical model. It is not a \emph{physical} validation, and it says nothing
about how well a homogeneous model represents the true gravity field of Ryugu.
That limitation is taken up in Section~\ref{sec:uniformdensity}.

The comparison is run on the full mesh, with no decimation, so that our facets
correspond one-to-one with the archived plates. The correspondence is confirmed
geometrically rather than assumed: matching by index order, the median distance
between our facet centroids and the archived plate centres is $0.04$~m, with a
maximum of $0.1$~m, against a body some $900$~m across. The two meshes are
therefore the same mesh in the same order, and a facet-by-facet difference is
meaningful.

Table~\ref{tab:ryugu} summarises the outcome, and
Figure~\ref{fig:ryuguval} shows the comparison facet by facet. The area-weighted mean slope
agrees to better than $0.01^{\circ}$, and the distribution of per-facet
differences is tight: the RMS difference is $0.48^{\circ}$ and the median
absolute difference $0.34^{\circ}$, with $49\,133$ of the $49\,152$ facets
($99.96\%$) agreeing to within $5^{\circ}$ and $49\,127$ ($99.95\%$) to within
$2^{\circ}$. The Pearson correlation across facets is $0.9989$ for slope and
$0.9898$ for geopotential.

The residual is not uniformly small, and it is worth being explicit about its
tail. Nineteen facets differ by more than $5^{\circ}$ and seven by more than
$10^{\circ}$, the largest by $24.45^{\circ}$. These are a vanishing fraction of
the surface---the $99.9$th percentile of the absolute difference is
$0.95^{\circ}$---but they are real, and Figure~\ref{fig:ryuguval}(b) is plotted
on a logarithmic count axis over the full residual range so that they remain
visible.

\begin{table}[t]
\centering
\small
\caption{Ryugu, facet-by-facet comparison against the archived Hayabusa2
gravitational products. Both fields are computed on the same $49\,152$-facet
shape model at $\dens=1200\ \mathrm{kg\,m^{-3}}$ and $P=7.63$~h.}
\label{tab:ryugu}
\begin{tabular}{lr}
\toprule
Quantity & Value \\
\midrule
Mean slope, this work                       & $15.06^{\circ}$ \\
Mean slope, archived products               & $15.06^{\circ}$ \\
Mean difference                             & $-0.00^{\circ}$ \\
RMS difference                              & $0.48^{\circ}$ \\
Median absolute difference                  & $0.34^{\circ}$ \\
$99.9$th percentile of $\lvert$difference$\rvert$ & $0.95^{\circ}$ \\
Maximum absolute difference                 & $24.45^{\circ}$ \\
Facets within $2^{\circ}$                   & $49\,127$ \ ($99.95\%$) \\
Facets within $5^{\circ}$                   & $49\,133$ \ ($99.96\%$) \\
Facets differing by more than $5^{\circ}$   & $19$ \\
Facets differing by more than $10^{\circ}$  & $7$ \\
Pearson $r$, slope                          & $0.9989$ \\
Pearson $r$, geopotential                   & $0.9898$ \\
Pearson $r$, slope on facets below $15^{\circ}$ & $0.9936$ \\
\bottomrule
\end{tabular}
\end{table}

Two features of the residual are worth noting. First, removing the mean bias
does not reduce the RMS difference, which stays at $0.48^{\circ}$: the
disagreement is scatter rather than offset, as expected if it arises from
differences in how the two implementations evaluate the field near the surface
rather than from a systematic error in either. Second, the correlation is just
as high on the flattest part of the surface---$0.9936$ restricted to facets
whose archived slope is below $15^{\circ}$---as it is overall. A sign or
convention error would degrade the flat facets preferentially, because there
the two competing contributions are most nearly balanced; that the flat facets
agree as well as the steep ones is a useful independent indication that the
conventions of Section~\ref{sec:signs} are correct.

The agreement also holds as a function of position on the body. Binning both
data sets by latitude gives the profile in Table~\ref{tab:ryugulat}: a maximum
near the equator, a minimum at intermediate latitude, and a rise toward the
poles. All four bands agree to better than $0.7^{\circ}$, and the two lower
bands to better than $0.2^{\circ}$. The equatorial maximum and mid-latitude
minimum are the signature of the equatorial ridge of a top-shaped body, and the
framework reproduces both without adjustment.

\begin{table}[t]
\centering
\small
\caption{Ryugu, area-weighted mean slope by latitude band, compared with the
archived products.}
\label{tab:ryugulat}
\begin{tabular}{lrrr}
\toprule
$\lvert\text{latitude}\rvert$ & This work & Archive & Difference \\
\midrule
$0$--$20^{\circ}$   & $15.12^{\circ}$ & $15.20^{\circ}$ & $-0.08^{\circ}$ \\
$20$--$40^{\circ}$  & $12.23^{\circ}$ & $12.40^{\circ}$ & $-0.17^{\circ}$ \\
$40$--$60^{\circ}$  & $15.66^{\circ}$ & $15.00^{\circ}$ & $+0.66^{\circ}$ \\
$60$--$90^{\circ}$  & $21.25^{\circ}$ & $20.80^{\circ}$ & $+0.45^{\circ}$ \\
\midrule
Global              & $15.06^{\circ}$ & $15.06^{\circ}$ & $-0.00^{\circ}$ \\
\bottomrule
\end{tabular}
\end{table}

\begin{figure}[tp]
\centering
\includegraphics[width=\textwidth]{fig1.pdf}
\caption{Per-facet validation of the computed slope field against the
gravitational products released with the Hayabusa2 shape analysis of Ryugu,
on the identical $49\,152$-facet model at $\dens = 1200\ \mathrm{kg\,m^{-3}}$
and $P = 7.63$~h. (a) Facet-by-facet comparison; the colour scale is
logarithmic in the number of facets per bin, and the dashed line is the $1{:}1$
relation. (b) Distribution of the per-facet difference. The agreement is
scatter rather than offset: removing the mean bias leaves the RMS difference
unchanged.}
\label{fig:ryuguval}
\end{figure}

\subsection{Bennu: global structure}
\label{sec:bennu}

For Bennu the OLA shape model \citep{barnouin2020,daly2020} is decimated to
$30\,561$ facets and evaluated at the measured bulk density
$\dens=1.194\ \mathrm{g\,cm^{-3}}$ and the rotation period $P=4.2961$~h. The
resulting area-weighted mean slope is $15.45^{\circ}$, with a median of
$14.94^{\circ}$; the distribution is concentrated between $10^{\circ}$ and
$20^{\circ}$ (which together hold $48.7\%$ of the surface area), only $3.9\%$
of facets exceed $30^{\circ}$, and just two facets out of $30\,561$ exceed
$90^{\circ}$. These figures sit within the range of global mean slopes reported
from OSIRIS-REx analyses \citep{scheeres2019,jawin2020}.

The latitudinal structure is more diagnostic than the global mean, because it
is set by the competition between the equatorial ridge and the centrifugal
term. The profile, given in Table~\ref{tab:bennulat}, rises from $10.45^{\circ}$
within $10^{\circ}$ of the equator to a broad maximum of $19.89^{\circ}$ at
$45$--$60^{\circ}$ latitude before falling to $17.83^{\circ}$ near the pole.
This reproduces the pattern reported from mission data---low slopes across the
equatorial region, a transition near $\pm 20^{\circ}$, and higher slopes at
mid to high latitude \citep{scheeres2019}---both in shape and, to within a
degree or two, in value; Figure~\ref{fig:bennuval}(a) shows the profile against
the published values.

\begin{table}[t]
\centering
\small
\caption{Bennu, area-weighted mean slope by latitude band at
$\dens = 1.194\ \mathrm{g\,cm^{-3}}$ and $\Nf = 30\,561$.}
\label{tab:bennulat}
\begin{tabular}{lrrr}
\toprule
$\lvert\text{latitude}\rvert$ & Mean & Median & Facets \\
\midrule
$0$--$10^{\circ}$   & $10.45^{\circ}$ & $9.70^{\circ}$  & $6202$ \\
$10$--$20^{\circ}$  & $12.02^{\circ}$ & $11.36^{\circ}$ & $5804$ \\
$20$--$30^{\circ}$  & $16.06^{\circ}$ & $15.48^{\circ}$ & $4949$ \\
$30$--$45^{\circ}$  & $19.10^{\circ}$ & $18.43^{\circ}$ & $5856$ \\
$45$--$60^{\circ}$  & $19.89^{\circ}$ & $19.32^{\circ}$ & $4230$ \\
$60$--$90^{\circ}$  & $17.83^{\circ}$ & $17.36^{\circ}$ & $3520$ \\
\bottomrule
\end{tabular}
\end{table}

The geopotential of Eq.~\eqref{eq:geopot} carries the same information in a
form directly relevant to regolith transport. Its minimum lies at the equator
and it increases toward the poles, so that loose material on the surface is
driven equatorward; this is the geopotential structure inferred for Bennu from
mission data and invoked to explain its equatorial accumulation
\citep{scheeres2019,jawin2020}. The framework recovers it from the shape model, the
assumed density and the specified rotation period.

\subsection{Bennu: response to the assumed density}
\label{sec:densitysweep}

A shape-based calculation must assume a bulk density, and it is worth
establishing that the framework responds to that assumption in the physically
correct way rather than merely producing a plausible number at one value. For a
rotating body the relevant control parameter is the ratio of centrifugal to
gravitational acceleration; since the centrifugal term is fixed by the spin and
the shape, raising the assumed density strengthens self-gravity and reduces
that ratio, which must lower the slope.

Because the field is linear in $\dens$ (Section~\ref{sec:solver}), the entire
sweep costs one solver evaluation. Table~\ref{tab:sweep} gives the result, and
Figure~\ref{fig:bennuval}(b) displays it. The
mean slope falls from $25.55^{\circ}$ at $\dens = 0.85\ \mathrm{g\,cm^{-3}}$ to
$15.45^{\circ}$ at the measured $1.194\ \mathrm{g\,cm^{-3}}$, while the
equatorial centrifugal-to-gravitational ratio falls from $0.751$ to $0.534$.
The comparison values reported from OSIRIS-REx analyses for the same exercise
are approximately $24^{\circ}$ at $0.85$ and $15^{\circ}$ at
$1.15\ \mathrm{g\,cm^{-3}}$ \citep{scheeres2019}; the framework returns
$25.55^{\circ}$ and $16.06^{\circ}$, agreeing to $1.5^{\circ}$ and
$1.1^{\circ}$ respectively. The agreement is therefore quantitative across the
range, not merely in trend, and at the measured density the equatorial and
polar means---$11.19^{\circ}$ and $17.81^{\circ}$---bracket the published
equatorial and high-latitude figures of roughly $12^{\circ}$ and $18^{\circ}$.

\begin{table}[t]
\centering
\small
\caption{Bennu, area-weighted mean slope as a function of the assumed bulk
density at $\Nf = 30\,561$. Here $\text{eq}$ denotes
$\lvert\text{lat}\rvert<20^{\circ}$, $\text{mid}$ denotes
$20$--$40^{\circ}$, and $\text{pol}$ denotes $>60^{\circ}$;
$\lVert\acf\rVert/\lVert\ngrav\rVert$ is the equatorial mean.}
\label{tab:sweep}
\begin{tabular}{rrrrrr}
\toprule
$\dens$ & Mean & eq & mid & pol & $\lVert\acf\rVert/\lVert\ngrav\rVert$ \\
\midrule
$0.850$ & $25.55^{\circ}$ & $23.07^{\circ}$ & $29.38^{\circ}$ & $21.28^{\circ}$ & $0.751$ \\
$1.000$ & $19.18^{\circ}$ & $14.57^{\circ}$ & $22.49^{\circ}$ & $19.43^{\circ}$ & $0.638$ \\
$1.150$ & $16.06^{\circ}$ & $11.61^{\circ}$ & $18.19^{\circ}$ & $18.12^{\circ}$ & $0.555$ \\
$1.194$ & $15.45^{\circ}$ & $11.19^{\circ}$ & $17.26^{\circ}$ & $17.81^{\circ}$ & $0.534$ \\
\bottomrule
\end{tabular}
\end{table}

\begin{figure}[tp]
\centering
\includegraphics[width=\textwidth]{fig2.pdf}
\caption{Bennu. (a) Area-weighted mean and median slope by latitude band at the
measured bulk density; shaded bands mark the equatorial and high-latitude values
reported from OSIRIS-REx analyses. The profile reproduces the observed pattern
of low equatorial slopes rising to a mid-latitude maximum. (b) Dependence of the
slope on the assumed bulk density, resolved by latitude zone. Stars mark the
published global means at $0.85$ and $1.15\ \mathrm{g\,cm^{-3}}$. The
equatorial zone responds most strongly, as expected for a centrifugal effect
that scales with distance from the spin axis.}
\label{fig:bennuval}
\end{figure}

The equatorial band responds most strongly, falling by nearly $12^{\circ}$
across the range while the polar band moves by only $3.5^{\circ}$. That is the
expected signature: the centrifugal term acts perpendicular to the spin axis
and so has its largest effect where the surface is farthest from that axis.
The framework therefore reproduces not only the magnitude of the
density dependence but its spatial distribution.

\subsection{Summary}
\label{sec:valsummary}

Taken together, the two tests bound the accuracy of the framework from
different directions. Ryugu establishes that, given identical inputs, the
implementation reproduces an independent Werner--Scheeres calculation facet by
facet, with an RMS difference of half a degree and $99.96\%$ of facets
agreeing to within $5^{\circ}$. Bennu establishes that, given a shape model, an
assumed density and a specified spin state, the framework recovers the
published global mean, the
latitudinal structure, the equatorial geopotential minimum, and the
quantitative dependence on assumed density.

Neither test says anything about how these numbers would change had a different
mesh been used, and that question is the subject of the next two sections.

\section{Resolution behaviour of the slope}
\label{sec:slopesens}

Section~\ref{sec:validation} established that the framework reproduces
independent results when the mesh is fixed. This section asks a different
question: how much of a reported slope is a property of the body, and how much
is a property of the mesh on which it was computed.

\subsection{The measurement}
\label{sec:slopemeas}

For each body a single delivered shape model is decimated to a sequence of five
target facet counts spanning roughly $2\,000$ facets to the full model, using
the vertex-clustering operator of Section~\ref{sec:operators}. Density and
rotation period are held fixed at the values of Table~\ref{tab:shapes}, so the
body, its mass, and its spin are identical across the sequence; only the
sampling of its surface changes. The area-weighted mean slope $\slopeavg$ of
Eq.~\eqref{eq:slopeavg} is evaluated on each mesh.

Five bodies are used: Ryugu, Eros, Itokawa, Gaspra, and Ida. Kleopatra is
excluded, because both available models are already coarse ($4\,092$ and
$2\,292$ facets) and saturate under decimation, leaving too few distinct levels
to characterise a trend.

\subsection{Results}
\label{sec:slopesensres}

Table~\ref{tab:slopeseries} gives the full sequences, and the upper panel of
Figure~\ref{fig:resolution} plots them as fractional departures from the
finest-mesh value. On every body without exception $\slopeavg$ increases
monotonically with facet count. The increase is
not marginal: over the tested range it amounts to $3.5^{\circ}$ on Ryugu,
$2.2^{\circ}$ on Eros and $2.0^{\circ}$ on Itokawa, on bodies whose mean slopes
are themselves only $11$--$17^{\circ}$.

\begin{table}[t]
\centering
\small
\caption{Area-weighted mean slope $\slopeavg$ as a function of facet count for
a single shape model of each body, decimated by vertex clustering. Density and
rotation period are fixed at the values of Table~\ref{tab:shapes}.}
\label{tab:slopeseries}
\begin{tabular}{lrr}
\toprule
Body & $\Nf$ & $\slopeavg$ \\
\midrule
Ryugu   & $3\,096$  & $12.174^{\circ}$ \\
        & $5\,780$  & $12.829^{\circ}$ \\
        & $11\,134$ & $13.688^{\circ}$ \\
        & $20\,014$ & $14.505^{\circ}$ \\
        & $42\,834$ & $15.700^{\circ}$ \\
\midrule
Eros    & $2\,097$  & $11.034^{\circ}$ \\
        & $6\,306$  & $11.691^{\circ}$ \\
        & $12\,155$ & $12.122^{\circ}$ \\
        & $21\,850$ & $12.705^{\circ}$ \\
        & $36\,030$ & $13.263^{\circ}$ \\
\midrule
Itokawa & $2\,183$  & $14.710^{\circ}$ \\
        & $6\,404$  & $15.167^{\circ}$ \\
        & $12\,272$ & $15.762^{\circ}$ \\
        & $21\,196$ & $16.304^{\circ}$ \\
        & $41\,369$ & $16.727^{\circ}$ \\
\midrule
Gaspra  & $2\,821$  & $12.413^{\circ}$ \\
        & $5\,240$  & $12.466^{\circ}$ \\
        & $12\,500$ & $12.564^{\circ}$ \\
        & $20\,773$ & $12.619^{\circ}$ \\
        & $32\,040$ & $12.641^{\circ}$ \\
\midrule
Ida     & $2\,141$  & $13.171^{\circ}$ \\
        & $5\,808$  & $13.560^{\circ}$ \\
        & $11\,804$ & $13.807^{\circ}$ \\
        & $20\,673$ & $13.918^{\circ}$ \\
        & $32\,040$ & $13.933^{\circ}$ \\
\bottomrule
\end{tabular}
\end{table}

The two diagnostics of Section~\ref{sec:diagnostics} summarise these sequences
and, as anticipated there, they do not tell the same story
(Table~\ref{tab:slopesens}; see also the left-hand bars of each panel of
Figure~\ref{fig:diagnostics}).

The range sensitivity $S[\slopeavg]$ spans an order of magnitude between
bodies, from $1.8\%$ on Gaspra to $25.6\%$ on Ryugu. A quarter of the reported
value on Ryugu, and nearly a fifth on Eros, is attributable to the choice of
mesh alone.

The convergence residual $C[\slopeavg]$ is the more troubling of the two. On
Gaspra and Ida the slope has effectively settled by the finest mesh, changing
by only $0.17\%$ and $0.11\%$ over the last refinement step. On Ryugu, Eros and
Itokawa it has not: the last step still moves the value by $8.24\%$, $4.39\%$
and $2.59\%$ respectively, with no sign of a plateau. For those bodies there is
no evidence, within the resolution available, that $\slopeavg$ tends to a limit
at all. A value quoted at the finest mesh is therefore not an estimate of a
converged quantity; it is the value at that mesh.

\begin{table}[t]
\centering
\small
\caption{Range sensitivity $S$ and convergence residual $C$ of the
area-weighted mean slope, from the sequences of Table~\ref{tab:slopeseries}.
$C$ is the fractional change over the final refinement step.}
\label{tab:slopesens}
\begin{tabular}{lrrr}
\toprule
Body & $\Nf$ range & $S[\slopeavg]$ & $C[\slopeavg]$ \\
\midrule
Ryugu   & $3\,096$--$42\,834$ & $25.6\%$ & $8.24\%$ \\
Eros    & $2\,097$--$36\,030$ & $18.3\%$ & $4.39\%$ \\
Itokawa & $2\,183$--$41\,369$ & $12.8\%$ & $2.59\%$ \\
Ida     & $2\,141$--$32\,040$ & $5.6\%$  & $0.11\%$ \\
Gaspra  & $2\,821$--$32\,040$ & $1.8\%$  & $0.17\%$ \\
\bottomrule
\end{tabular}
\end{table}

\subsection{The origin of the trend}
\label{sec:slopeorigin}

That the slope rises with refinement, always in the same direction, points to a
specific mechanism. Decimation does not remove relief from the surface
(Section~\ref{sec:operators}); a coarse facet spans the same topography that
several fine facets spanned before, and its normal is close to the
area-weighted average of theirs. Averaging unit normals over a rough patch
yields a direction more nearly aligned with the mean surface than any individual
normal, so a coarse mesh systematically under-represents local departures from
the mean figure. Refining the mesh restores them, and the mean slope rises
accordingly.

If that account is right, then removing the relief at fixed facet count should
lower the slope by a comparable amount. The Laplacian smoothing operator of
Eq.~\eqref{eq:laplace} tests this directly: it suppresses short-wavelength
relief while leaving $\Nf$, the global figure, the mass and the spin unchanged.
Table~\ref{tab:smooth} gives the outcome for Itokawa and Eros.

\begin{table}[t]
\centering
\small
\caption{Effect of Laplacian smoothing at fixed facet count
($\Nf = 34\,234$ for Itokawa, $36\,030$ for Eros). The last column is the
fraction of facets steeper than $30^{\circ}$.}
\label{tab:smooth}
\begin{tabular}{lrrrr}
\toprule
Body & Iterations & $\slopeavg$ & Median & $>30^{\circ}$ \\
\midrule
Itokawa & $0$  & $15.81^{\circ}$ & $12.34^{\circ}$ & $6.5\%$ \\
        & $1$  & $14.82^{\circ}$ & $11.53^{\circ}$ & $4.9\%$ \\
        & $3$  & $14.00^{\circ}$ & $10.98^{\circ}$ & $4.3\%$ \\
        & $5$  & $13.56^{\circ}$ & $10.77^{\circ}$ & $3.9\%$ \\
        & $10$ & $12.90^{\circ}$ & $10.51^{\circ}$ & $3.4\%$ \\
\midrule
Eros    & $0$  & $13.65^{\circ}$ & $12.05^{\circ}$ & $3.5\%$ \\
        & $1$  & $12.23^{\circ}$ & $10.87^{\circ}$ & $1.6\%$ \\
        & $3$  & $11.35^{\circ}$ & $10.11^{\circ}$ & $0.8\%$ \\
        & $5$  & $10.92^{\circ}$ & $9.74^{\circ}$  & $0.6\%$ \\
        & $10$ & $10.30^{\circ}$ & $9.25^{\circ}$  & $0.2\%$ \\
\bottomrule
\end{tabular}
\end{table}

Smoothing lowers the mean slope by $2.9^{\circ}$ on Itokawa and $3.4^{\circ}$
on Eros, comparable to the $2.0^{\circ}$ and $2.2^{\circ}$ that refinement adds
over the tested range, and it acts most strongly on the steep tail: the fraction
of facets above $30^{\circ}$ falls by roughly half on Itokawa and by more than
an order of magnitude on Eros. The two operators move $\slopeavg$ in opposite
directions for the same underlying reason, and the mechanism is confirmed. What
the mesh resolution controls is how much of the surface relief the slope field
is able to see.

This also explains why the effect is a property of the shape model rather than
of the asteroid. Gaspra and Ida, whose models derive from limited flyby imaging
and carry little sub-facet relief to recover, show the smallest sensitivities
and have converged by the finest mesh. Ryugu, whose model is a dense
photogrammetric reconstruction rich in unresolved detail, shows the largest.
$S[\slopeavg]$ is therefore better read as a measure of how much relief a shape
model still contains below its nominal resolution than as a characteristic of
the body itself.

\subsection{Consequences for reporting}
\label{sec:slopereport}

Three consequences follow, and they are practical rather than conceptual.

First, a mean slope quoted without its facet count is not reproducible. On three
of the five bodies here the same shape model yields values differing by more
than $2^{\circ}$ depending on a choice that is rarely stated.

Second, mean slopes computed on different shape models of the same body are not
comparable, and neither are values from different studies unless both
resolutions are known and similar. Section~\ref{sec:crossversion} quantifies
this for a case where two independent models of one asteroid are available.

Third, and least comfortably, on the bodies with the richest shape models no
value within the tested range can be quoted as converged. Reporting $\slopeavg$ at the finest available
mesh is defensible only if the mesh is stated alongside it, and any comparison
must then be made at matched resolution. The next section shows that a closely
related quantity does not share this behaviour.

\section{Resolution robustness of the geopotential dispersion}
\label{sec:varrobust}

The difficulty identified in Section~\ref{sec:slopesens} is not that the slope
is noisy---it is smooth and monotonic in the mesh---but that it does not
approach a limit. This section shows that the normalized geopotential dispersion
$\potvarn$ of Eq.~\eqref{eq:potvarn}, computed from the same fields on the same
meshes, behaves differently: it too rises with refinement, but it saturates.

\subsection{The measurement}
\label{sec:varmeas}

The measurement is identical to that of Section~\ref{sec:slopemeas} and is
performed on the same decimation sequences, from the same evaluations. No
additional solver calls are involved: $\slopeavg$ and $\potvarn$ are two
summaries of one computed field, so any difference between them is a difference
between the diagnostics, not between the calculations that produced them.

\subsection{Results}
\label{sec:varres}

Table~\ref{tab:varseries} gives the sequences, and Figure~\ref{fig:resolution}
plots them directly beneath the corresponding slope curves, on identical axes.
As with the slope, $\potvarn$ increases monotonically with facet count on every
body. Unlike the slope, it
flattens: on Itokawa and Gaspra the last two meshes return the same value to
five significant figures, and on Ida they differ in the fifth.

\begin{table}[t]
\centering
\small
\caption{Normalized geopotential dispersion $\potvarn$ as a function of facet
count, on the same decimation sequences as Table~\ref{tab:slopeseries}.}
\label{tab:varseries}
\begin{tabular}{lrr}
\toprule
Body & $\Nf$ & $\potvarn$ \\
\midrule
Ryugu   & $3\,096$  & $0.03153$ \\
        & $5\,780$  & $0.03233$ \\
        & $11\,134$ & $0.03295$ \\
        & $20\,014$ & $0.03327$ \\
        & $42\,834$ & $0.03347$ \\
\midrule
Eros    & $2\,097$  & $0.03241$ \\
        & $6\,306$  & $0.03345$ \\
        & $12\,155$ & $0.03384$ \\
        & $21\,850$ & $0.03403$ \\
        & $36\,030$ & $0.03413$ \\
\midrule
Itokawa & $2\,183$  & $0.08527$ \\
        & $6\,404$  & $0.08577$ \\
        & $12\,272$ & $0.08598$ \\
        & $21\,196$ & $0.08608$ \\
        & $41\,369$ & $0.08608$ \\
\midrule
Gaspra  & $2\,821$  & $0.05929$ \\
        & $5\,240$  & $0.05980$ \\
        & $12\,500$ & $0.06009$ \\
        & $20\,773$ & $0.06017$ \\
        & $32\,040$ & $0.06017$ \\
\midrule
Ida     & $2\,141$  & $0.03490$ \\
        & $5\,808$  & $0.03572$ \\
        & $11\,804$ & $0.03626$ \\
        & $20\,673$ & $0.03641$ \\
        & $32\,040$ & $0.03643$ \\
\bottomrule
\end{tabular}
\end{table}

\begin{figure}[tp]
\centering
\includegraphics[width=\textwidth]{fig3.pdf}
\caption{Behaviour of the two diagnostics under mesh refinement, expressed as
the fractional deviation from the value obtained on the finest mesh of each
sequence. (a) Area-weighted mean slope. (b) Normalized geopotential dispersion,
plotted on \emph{identical} axes. Both quantities are computed from the same
field on the same meshes, so the contrast is a property of the diagnostics
rather than of the calculation. The slope curves spread by up to a quarter of
their own value and are still moving at the finest meshes; the dispersion curves
are flat on this scale.}
\label{fig:resolution}
\end{figure}

\subsection{The two diagnostics compared}
\label{sec:varcompare}

Table~\ref{tab:varsens} and Figure~\ref{fig:diagnostics} place the two
diagnostics side by side. We report both
measures of Section~\ref{sec:diagnostics}, because they support claims of
different strength and it would be misleading to quote only the more favourable
one.

By the range sensitivity $S$, the geopotential dispersion is the more stable
diagnostic on every body, but the margin varies considerably: it is a factor of
$13.6$ on Itokawa, $4.3$ on Ryugu and $3.6$ on Eros, but only $1.3$ on Ida and
$1.2$ on Gaspra. The last two figures should not be read as a weakness of the
dispersion. Gaspra and Ida are precisely the bodies whose slope was already
stable ($S[\slopeavg]=1.8\%$ and $5.6\%$), so there was little for a better
diagnostic to improve upon; the ratio is uninformative where the denominator is
already small. The honest summary of the range metric is that $\potvarn$ varies
by $0.9\%$ to $5.9\%$ where $\slopeavg$ varies by $1.8\%$ to $25.6\%$, and that
the advantage is largest exactly where it matters most.

\begin{table}[t]
\centering
\small
\caption{Range sensitivity $S$ and convergence residual $C$ for the two
diagnostics, on identical meshes. A dash indicates that the diagnostic did not
change over the final refinement step at the precision retained.}
\label{tab:varsens}
\begin{tabular}{lrrrr}
\toprule
& \multicolumn{2}{c}{$S$} & \multicolumn{2}{c}{$C$} \\
\cmidrule(lr){2-3}\cmidrule(lr){4-5}
Body & $\slopeavg$ & $\potvarn$ & $\slopeavg$ & $\potvarn$ \\
\midrule
Ryugu   & $25.6\%$ & $5.9\%$ & $8.24\%$ & $0.60\%$ \\
Eros    & $18.3\%$ & $5.1\%$ & $4.39\%$ & $0.29\%$ \\
Itokawa & $12.8\%$ & $0.9\%$ & $2.59\%$ & $-$      \\
Ida     & $5.6\%$  & $4.3\%$ & $0.11\%$ & $0.05\%$ \\
Gaspra  & $1.8\%$  & $1.5\%$ & $0.17\%$ & $-$      \\
\bottomrule
\end{tabular}
\end{table}

\begin{figure}[tp]
\centering
\includegraphics[width=\textwidth]{fig4.pdf}
\caption{Resolution diagnostics for the two quantities.
(a) Range sensitivity $S$, the peak-to-peak variation over the full mesh range
normalized by the mean. (b) Convergence residual $C$, the fractional change over
the final refinement step. The geopotential dispersion is the more stable diagnostic
on every body by both measures, but the two measures support claims of different
strength: the range advantage is modest on bodies whose slope is already stable,
whereas the convergence residual separates the two cleanly---the dispersion has
reached a plateau on every body, and the slope has not on Ryugu, Eros or
Itokawa.}
\label{fig:diagnostics}
\end{figure}

The convergence residual $C$ draws a sharper and, we argue, more meaningful
distinction. On the three bodies where the slope has not settled---Ryugu, Eros
and Itokawa---the geopotential dispersion has: it changes by $0.60\%$, $0.29\%$ and
nothing at all over the same final refinement step on which the slope still
moves by $8.24\%$, $4.39\%$ and $2.59\%$. On the two bodies where the slope has
settled, the dispersion has settled at least as firmly. There is no body in the
set on which the dispersion is the less converged of the two.

The difference between the two rows of Table~\ref{tab:varsens} is therefore not
merely quantitative. Within the resolutions available, $\potvarn$ appears to
possess a limit that the meshes are approaching, and $\slopeavg$ on the richest
models does not.

\subsection{Why the two diagnostics differ in resolution sensitivity}
\label{sec:whyconverge}

The two diagnostics are computed from the same field, but they depend on the
mesh in fundamentally different ways.

The geopotential of Eq.~\eqref{eq:geopot} is an integral over the body's
volume. Refining the mesh changes the polyhedron's volume and its mass
distribution only slightly, and those changes are what the potential responds
to; a facet subdivided into four contributes very nearly what it contributed
before, because the mass it bounds is nearly unchanged. The dispersion inherits
this behaviour, and both $\potvar$ and $\Ubar$ converge as the polyhedron
converges to the body it approximates.

The slope of Eq.~\eqref{eq:slope} depends on the facet normal $\nhat_{i}$, a
property of the surface at the scale of a single facet. Subdividing a facet
does not leave the normals nearly unchanged: it replaces one direction by
several that differ from it and from each other by whatever relief exists at the
finer scale. There is no reason for that sequence of directions to settle,
because a rough surface continues to present new orientations at every scale at
which it is resolved. This is the familiar situation in which an enclosed volume
converges under refinement while a surface property need not, and it is why the
distinction of Section~\ref{sec:diagnostics} between $S$ and $C$ was worth
drawing: the slope can be made to look stable by choosing a narrow enough range
of meshes, but it does not thereby acquire a limit.

The smoothing experiment of Table~\ref{tab:smooth} supports the same reading
from the other direction. Suppressing the fine relief---removing the very
structure that refinement keeps discovering---lowers the slope substantially,
by $2.9^{\circ}$ on Itokawa and $3.4^{\circ}$ on Eros. It is that structure,
not the body's overall figure, that the slope keeps responding to.

\subsection{What the geopotential dispersion measures}
\label{sec:varmeaning}

Beyond its numerical stability, $\potvarn$ has a direct physical reading. It is
the area-weighted spread of the geopotential over the surface, normalized by
its mean, and therefore measures how far the surface departs from an
equipotential. A body whose surface has fully relaxed into an equipotential
figure would have $\potvarn = 0$; a body far from equilibrium has a large value.

This is what makes it useful beyond bookkeeping. The values in
Table~\ref{tab:varseries} separate the bodies more cleanly than the slope does:
Itokawa at $0.086$ and Gaspra at $0.060$ sit well above Ryugu, Eros and Ida,
which cluster between $0.033$ and $0.036$. Because the quantity is
dimensionless and resolution-stable, such a comparison is meaningful in a way
that a comparison of mean slopes across differently meshed models is not.

The same property underlies its use as an estimator. \citet{kanamaru2019}
minimized a closely related quantity over trial interior density distributions
to infer the internal structure of Itokawa, on the reasoning that a surface
shaped by long-term relaxation should lie close to an equipotential of the true
density field. Any such inversion requires a diagnostic whose variation with
the assumed density is not swamped by its variation with the mesh---which is
exactly the property established here, and which the slope does not possess.
We return to this in Section~\ref{sec:discussion}.

\subsection{Recommendation}
\label{sec:varrecommend}

The results of this section and the last suggest a simple practice. Where a
resolution-independent statement about a body's surface gravity field is
required---for comparison between bodies, between studies, or between shape
models---the normalized geopotential dispersion should be reported, since within
the range tested it converges. The mean slope remains the more directly
interpretable quantity, and nothing here argues against computing it; but it
should be quoted together with the facet count at which it was obtained, and
compared only at matched resolution.

The next section tests both diagnostics under the harder condition that the
shape models themselves differ.

\section{Comparison across shape-model versions}
\label{sec:crossversion}

Sections~\ref{sec:slopesens} and~\ref{sec:varrobust} varied the mesh while
holding the shape model fixed. This section varies the shape model itself. The
question is practical: most asteroids will never be visited, and their shapes
are known only from ground-based radar or lightcurve inversion. Which surface
gravity quantities, if any, survive that substitution?

\subsection{The controlled case: Itokawa}
\label{sec:itokawaxv}

Itokawa is the rare object for which two independent shape models of comparable
extent exist: the photogrammetric model built from Hayabusa imaging, with
$49\,152$ facets, and a radar-derived model with $3\,688$. Their bounding boxes
agree to $6.5\%$, $3.9\%$ and $13.3\%$ along the three principal axes, so the
two models describe a body of the same size and gross figure; what differs is
the detail they resolve. Both are evaluated at the same density and rotation
period, so the comparison isolates the effect of shape-model provenance.

Because Section~\ref{sec:slopesens} showed that resolution alone moves the
slope, the two models must be compared at matched resolution or the test is
meaningless. The radar model has $3\,688$ facets; the corresponding
spacecraft-model values of the bulk descriptors are obtained by linear
interpolation of the sequences of Tables~\ref{tab:slopeseries}
and~\ref{tab:varseries} between the $2\,183$- and $6\,404$-facet meshes. The
distribution statistics cannot be interpolated in the same way and are taken
from the nearest available direct runs, at $4\,857$ facets for the spacecraft
model against $3\,688$ for the radar model; the comparison of the tails is
therefore at approximately rather than exactly matched resolution, a caveat we
return to below. Table~\ref{tab:itokawaxv} gives the result, and
Figure~\ref{fig:crossversion}(a) shows the two slope distributions.

\begin{table}[t]
\centering
\small
\caption{Itokawa: spacecraft-derived and radar-derived shape models compared at
$\Nf \simeq 3\,688$, at $\dens = 1.95\ \mathrm{g\,cm^{-3}}$ and
$P = 12.132$~h. Spacecraft-model values for the bulk descriptors are
interpolated to the radar model's facet count; distribution statistics are from
direct runs at $4\,857$ and $3\,688$ facets respectively.}
\label{tab:itokawaxv}
\begin{tabular}{lrrr}
\toprule
Quantity & Spacecraft & Radar & Difference \\
\midrule
$\slopeavg$ (area-weighted)   & $14.87^{\circ}$ & $14.75^{\circ}$ & $-0.8\%$  \\
$\potvarn$                    & $0.08545$       & $0.07874$       & $-7.9\%$  \\
\midrule
Mean slope (unweighted)       & $15.91^{\circ}$ & $13.40^{\circ}$ & $-15.8\%$ \\
Maximum slope                 & $149.45^{\circ}$& $41.32^{\circ}$ & --- \\
Facets above $30^{\circ}$     & $4.8\%$         & $1.7\%$         & --- \\
Facets above $90^{\circ}$     & $2.29\%$        & $0.00\%$        & --- \\
\bottomrule
\end{tabular}
\end{table}

\begin{figure}[tp]
\centering
\includegraphics[width=\textwidth]{fig5.pdf}
\caption{Shape-model detail and the slope distribution. (a) Slope distribution
of Itokawa from the spacecraft-derived ($N_f = 4\,857$) and radar-derived
($N_f = 3\,688$) shape models. The bulk descriptors, compared at a common facet
count by interpolation, agree closely---their area-weighted mean slopes differ
by $0.8\%$---but the radar model contains no steep tail.
(b) Effect of Laplacian smoothing at fixed facet count. Ten iterations bring the
spacecraft model of Itokawa to $12.90^{\circ}$, within $3.7\%$ of the value
obtained independently from the radar model (dashed line), supporting the
interpretation that the difference between the two models is dominated by
facet-scale relief.}
\label{fig:crossversion}
\end{figure}

The two halves of that table point in opposite directions, and the contrast is
the result of this section.

The bulk descriptors, compared at the same facet count, agree well: the
area-weighted mean slope differs by less than one per cent and the normalized
geopotential dispersion by under eight per cent. A ground-based shape model, carrying an order of magnitude fewer facets
and no knowledge of the surface below the radar resolution, reproduces both
global summaries of the gravity field to within the accuracy with which the
assumed bulk density is usually known.

The distribution tails do not agree at all. The radar model contains no facet
steeper than $41^{\circ}$ and places $1.7\%$ of the surface above $30^{\circ}$;
the spacecraft model reaches $149^{\circ}$, places $4.8\%$ above $30^{\circ}$,
and assigns $2.29\%$ of its facets slopes exceeding $90^{\circ}$---overhanging
surfaces that the radar model does not represent at all. Nothing in the radar
model's global agreement gives any warning that its steep tail is absent.

The unweighted mean slope sits between the two behaviours, differing by
$15.8\%$ where the area-weighted mean differs by $0.8\%$. The reason is that
the spacecraft model's steep facets are also its smallest: they occupy a minor
fraction of the surface area but a substantial fraction of the facet count, so
they dominate an unweighted average and are correctly suppressed in an
area-weighted one. This is a further argument for the area weighting adopted in
Eq.~\eqref{eq:slopeavg}, and a caution that the two averages should not be used
interchangeably when models of different resolution are involved.

\subsection{Consistency with the smoothing experiment}
\label{sec:xvsmooth}

Section~\ref{sec:slopeorigin} attributed the resolution dependence of the slope
to relief that coarse meshes fail to resolve. If the difference between the two
Itokawa models is the same effect---the radar model simply never having
recorded that relief---then artificially removing the relief from the
spacecraft model should bring it toward the radar model's value.

It does. Ten Laplacian iterations applied to the spacecraft model
(Table~\ref{tab:smooth}) reduce its mean slope to $12.90^{\circ}$, within
$3.7\%$ of the $13.40^{\circ}$ obtained from the independent radar model, as
Figure~\ref{fig:crossversion}(b) shows. The
agreement is not exact and should not be over-read: smoothing and coarse
observation are different operations, and the residual difference is comparable
to what remains between the models generally. But that a purely internal
operation on one model reproduces, to within a few per cent, the value obtained
from an independently constructed model is a useful consistency check on the
interpretation.

\subsection{A cautionary counterexample: Kleopatra}
\label{sec:kleopatraxv}

Two shape models are also available for Kleopatra, and applying the same
comparison to them illustrates a prerequisite that is easy to overlook.

Compared at $\Nf \simeq 2\,292$, the two models disagree by $39.9\%$ in
area-weighted mean slope ($19.90^{\circ}$ against $11.96^{\circ}$) and by
$36.1\%$ in geopotential dispersion ($0.03683$ against $0.02354$). Taken at face
value this would overturn the Itokawa result.

It cannot be taken at face value. The two models do not describe the same body
at different levels of detail: their bounding boxes differ in scale. The
alternative model has a long axis of $276.3$~km against roughly $217$~km for
the radar model, while the intermediate and short axes agree to a few per cent,
so the two differ in aspect ratio ($3.6{:}1$ against $2.7{:}1$) as well as in
overall size. At fixed assumed density a longer body contains more mass, its
self-gravity is stronger relative to the fixed centrifugal term, and its slopes
are correspondingly lower---which is the direction of the observed
discrepancy. The $36$--$40\%$ differences therefore conflate scale, elongation
and surface detail, and cannot be attributed to shape-model provenance.

We report the comparison rather than discarding it, because the lesson is
useful. Before two shape models of one body are compared, their bounding boxes
should be checked against each other and against the published dimensions; this
is among the diagnostics listed in Section~\ref{sec:shapes} for exactly this
reason. Kleopatra fails that check and Itokawa passes it, which is why the
conclusions of this section rest on Itokawa alone.

\subsection{Implications for bodies known only from the ground}
\label{sec:xvimplications}

For the large majority of small bodies, a ground-based shape model is the only
one that will ever exist. The Itokawa comparison suggests a division of labour
for such objects.

In the one controlled case available to us, global descriptors of the surface
gravity field are largely preserved: the area-weighted mean slope is recovered
to within one per cent and the normalized geopotential dispersion to within
eight per cent, the latter comparable to the uncertainty contributed by the
assumed bulk density. Where
the goal is to characterise a population, to rank bodies by how far their
surfaces sit from an equipotential, or to provide first-order context for
mission planning, such models are adequate.

Anything that depends on the tail of the slope distribution is not recoverable.
Identifying the steepest terrain, estimating the fraction of surface above the
angle of repose, or locating overhangs requires relief that a ground-based
model does not contain, and the resulting figures will be systematically and
silently optimistic. A coarse model does not merely give a noisier answer to
these questions; it gives a confidently wrong one.

Demonstrated as it is for a single object, this comparison is not a general
calibration. It does, however, provide a first quantitative reference point for
the size of the discrepancy to be expected---close agreement in the bulk
descriptors, order-unity disagreement in the tails---and the same test can be
repeated wherever a second, independent shape model of a body becomes
available.

\section{Discussion}
\label{sec:discussion}

\subsection{What has been established}
\label{sec:established}

Three results stand on the evidence presented. First, the framework computes
the surface gravity field correctly: given identical inputs it reproduces an
independent Werner--Scheeres calculation facet by facet, with an RMS difference
of half a degree and $99.96\%$ of facets agreeing to within five degrees
(Section~\ref{sec:ryugu}), and given only a shape model and an assumed density
it recovers the published global mean, latitudinal structure, geopotential
pattern and density dependence for Bennu (Sections~\ref{sec:bennu}
and~\ref{sec:densitysweep}).

Second, the area-weighted mean slope retains a substantial mesh dependence over
the tested refinement range. On the three bodies with the most detailed shape
models it had not converged at the finest mesh available, still moving by
$2.6$--$8.2\%$ over the last refinement step (Section~\ref{sec:slopesens}).
This is a property of the diagnostic, and the smoothing experiment identifies
its cause as relief that coarse meshes fail to resolve.

Third, the normalized geopotential dispersion computed from the same field on
the same meshes shows clear practical convergence, changing by at most $0.6\%$
over the same final step
and by nothing at all on two bodies (Section~\ref{sec:varrobust}). The
distinction follows from what the two quantities depend on: a volume integral
of the mass distribution in one case, the orientation of an individual facet in
the other.

None of this makes the slope a less useful quantity---it remains the direct
measure of whether material will move---but it does mean that a slope value is
a statement about a particular mesh, and must be reported as one.

\subsection{The assumed density dominates the error budget}
\label{sec:densitydominates}

It would be a mistake to leave the impression that mesh resolution is the
principal difficulty facing shape-based surface gravity. For any body whose
mass has not been measured, the assumed bulk density is a far larger source of
uncertainty.

The Bennu sweep of Table~\ref{tab:sweep} makes the comparison concrete. Varying
the assumed density across the range that would be plausible for a carbonaceous
rubble pile in the absence of a mass determination, from $0.85$ to
$1.194\ \mathrm{g\,cm^{-3}}$, moves the mean slope by $10.1^{\circ}$. Mesh
resolution, on the bodies where we measured it, moves it by $2$--$3.5^{\circ}$.
The density assumption is thus the dominant term by a factor of three or more,
and it acts as a systematic shift rather than as scatter.

This ordering has a practical consequence. For a body known only from the
ground, refining the shape model buys less than constraining the mass would,
and the two are not substitutes. It also reinforces the argument of
Section~\ref{sec:xvimplications}: in the one controlled case available, a
coarse ground-based model reproduced the mean slope to within one per cent and
the geopotential dispersion to within eight per cent, while the density
assumption contributes considerably more, so the shape model is rarely the
limiting factor for global quantities.

\subsection{The uniform-density assumption}
\label{sec:uniformdensity}

Every field in this paper assumes a homogeneous interior. That assumption is
exact for none of the bodies considered, and it is worth being explicit about
where it bites and where it does not.

It does not affect the resolution results of Sections~\ref{sec:slopesens}
and~\ref{sec:varrobust}. Those hold density fixed and vary only the mesh, so
whatever error the homogeneity assumption introduces is common to every point
of a sequence and cancels from $S$ and $C$.

It does affect the absolute fields, and Itokawa is the clearest case. Its head
and body are inferred to differ in density by roughly a quarter, with the head
the denser of the two \citep{kanamaru2019}, a conclusion supported
independently by the YORP spin-up analysis of \citet{lowry2014}. A homogeneous
calculation must therefore misplace the local gravity direction near the neck,
where the two lobes contribute comparably. The comparison with the Ryugu
archive in Section~\ref{sec:ryugu} does not test this, since that reference is
itself computed for a uniform interior; it validates the solver, not the
physical assumption.

For bodies whose interiors are close to homogeneous---which the near-coincidence
of the centre of figure and centre of mass suggests for Eros---the assumption is
mild. For contact binaries and bilobed objects it is not, and slopes computed
near the junction should be treated with corresponding caution.

\subsection{Toward interior structure}
\label{sec:interior}

The limitation just described is also an opportunity, and it motivates the
companion work.

If a surface has been shaped over long periods by downslope transport, it should
tend toward an equipotential of the true gravity field, including whatever
density inhomogeneity that field contains. A trial interior density
distribution that is closer to the truth should therefore produce a smaller
spread of geopotential over the surface. This is the reasoning by which
\citet{kanamaru2019} inferred a denser head for Itokawa, and it converts the
geopotential dispersion from a descriptive statistic into an estimator.

Such an inversion places a specific demand on the diagnostic used: its variation
with the assumed density distribution must not be swamped by its variation with
the mesh. Section~\ref{sec:varrobust} establishes that $\potvarn$ satisfies
this, and Section~\ref{sec:slopesens} that the mean slope does not. That is the
main reason the distinction developed in this paper matters beyond questions of
reporting practice.

Two further ingredients are already in place. The linearity of the field in
density (Section~\ref{sec:solver}) means that a body partitioned into regions
of distinct density can be evaluated by superposing unit-density solutions, so
that scanning a density model costs one solver call per region rather than one
per trial. And Paper~I's bound $GM/\RA^{2}\le\gmean$ \citep{paper1} holds for
any interior mass distribution, providing a constraint that survives the
homogeneity assumption. Developing these into a working estimator, validating
it against the independently constrained cases, and establishing what can and
cannot be recovered from surface data alone is reserved for future work and is
not attempted here.

\subsection{On verification practice}
\label{sec:verifpractice}

One methodological point deserves emphasis beyond its role in this work.

The natural first test of a surface gravity code is a homogeneous sphere, which
must return zero slope everywhere. That test is necessary but weak: a
non-rotating sphere returns zero slope for either sign of the centrifugal term,
and so cannot detect an error in the rotating frame---which is precisely the
part of the calculation that distinguishes a small-body code from a
straightforward Newtonian one. The uniformly rotating sphere of
Eq.~\eqref{eq:rotsphere} does discriminate, and sharply: at
$\omegarot^{2}R/g = 0.5$ the two signs differ by more than seven degrees at
mid-latitude.

The cost of the stronger test is negligible---it is a closed-form expression and
a few lines of code, both provided with this work---and we suggest it be adopted
routinely. The same applies to checking the potential against the closed-form
ellipsoid solution of Eq.~\eqref{eq:ellipsoid}, which exercises the
gravitational term with the rotation switched off. Between them the two tests
cover every term in Eqs.~\eqref{eq:Ug}--\eqref{eq:geopot}, and neither depends
on any published data set.

\subsection{Limitations}
\label{sec:limitations}

The scope of what has been demonstrated should be stated plainly.

The per-facet validation rests on a single body. Ryugu is the only object for
which per-facet gravitational products computed on a publicly available shape
model at a stated density and spin are distributed, and it is therefore the only
case in which the implementation could be checked below the level of aggregate
statistics. That test is stringent, but it is one test.

The cross-version comparison rests on a single body as well. Of the two
asteroids in our set with independent shape models, only Itokawa passes the
prerequisite that the models describe a body of the same size
(Section~\ref{sec:kleopatraxv}). Whether the division between recoverable bulk
descriptors and unrecoverable distribution tails holds generally is not
established by one case.

The resolution study covers five bodies and a range of roughly $2\,000$ to
$42\,000$ facets. Whether the geopotential dispersion continues to hold its plateau
at $10^{5}$ or $10^{6}$ facets is untested, and the highest-resolution models
now available for Bennu and Ryugu would allow that range to be extended.

Finally, the slope of Eq.~\eqref{eq:slope} is one of several definitions in use.
It is the definition used by the Ryugu reference products, which is what makes
the comparison of Section~\ref{sec:ryugu} meaningful. Studies that define slope
relative to a geopotential-based reference surface rather than to the local
facet normal will obtain different values, particularly in equatorial regions of
rapidly rotating bodies, and the numbers here should not be compared with those
without accounting for the difference.

\section{Conclusions}
\label{sec:conclusions}

This paper extends the shape-based framework of Papers~I--III from the scalar
mean surface gravity to the spatially resolved surface fields---effective
gravity, gravitational slope, and geopotential---computed facet by facet from a
polyhedral shape model and an assumed uniform density. The main conclusions are
as follows.

\begin{enumerate}

\item \emph{The implementation reproduces an independent calculation facet by
facet.} Compared against the gravitational products released with the Hayabusa2
shape analysis of Ryugu, on the identical $49\,152$-facet model at the same
density and rotation period, the area-weighted mean slope agrees to better than
$0.01^{\circ}$, the RMS per-facet difference is $0.48^{\circ}$, $99.96\%$ of
facets agree to within $5^{\circ}$ and $99.95\%$ to within $2^{\circ}$, and the
Pearson
correlations are $0.9989$ in slope and $0.9898$ in geopotential. The agreement
holds band by band in latitude to better than $0.7^{\circ}$.

\item \emph{From a shape model, an assumed density and a specified spin state
the framework recovers the published surface gravity structure of Bennu.} The global mean slope is
$15.45^{\circ}$; the latitudinal profile rises from $10.45^{\circ}$ near the
equator to $19.89^{\circ}$ at mid-latitude; the geopotential minimum lies at the
equator, as inferred from mission data. Sweeping the assumed bulk density
reproduces the published slope--density relation quantitatively, returning
$25.55^{\circ}$ and $16.06^{\circ}$ at $0.85$ and $1.15\ \mathrm{g\,cm^{-3}}$
against roughly $24^{\circ}$ and $15^{\circ}$ reported from OSIRIS-REx analyses.

\item \emph{The area-weighted mean slope did not converge within the tested
mesh range on shape models that carry substantial unresolved relief.} Across
meshes from about $2\,000$ to $42\,000$ facets it varies by $1.8\%$ to $25.6\%$
depending on the body, and on three of five bodies it was still moving by
$2.6$--$8.2\%$ over the final refinement step, with no plateau. Laplacian
smoothing at fixed facet count lowers the slope by an amount comparable to what
refinement adds, identifying surface relief as the cause. A mean slope is
therefore a statement about a particular mesh and is not reproducible without
its facet count.

\item \emph{The normalized geopotential dispersion, computed from the same
field on the same meshes, shows clear practical convergence over the same
range.} It varies by $0.9\%$ to $5.9\%$ across the
same range and changes by at most $0.60\%$ over the final refinement step---by
nothing at all on two bodies---and is the more stable of the two diagnostics on
every body tested. The distinction follows from what each quantity depends on:
the geopotential is an integral over the mass distribution and inherits the
convergence of the polyhedron's volume, whereas the slope depends on the
orientation of an individual facet, which a rough surface continues to change
at every scale at which it is resolved.

\item \emph{In the controlled Itokawa comparison the ground-based model
preserves the global gravity descriptors but not the local slope statistics.}
Comparing the spacecraft-derived and radar-derived models at approximately
matched resolution, the area-weighted mean slope agrees to $0.8\%$ and the
geopotential dispersion to $7.9\%$, while the radar model contains no facet
steeper than $41^{\circ}$ against $149^{\circ}$ for the spacecraft model, and
places $1.7\%$ rather than $4.8\%$ of the surface above $30^{\circ}$. On the
evidence of this single controlled case, population-level characterisation from
ground-based shapes appears defensible, whereas identification of steep terrain,
overhangs, or the fraction of surface above the angle of repose does not.

\item \emph{A rotating test case is necessary to verify a small-body gravity
code.} A non-rotating sphere returns zero slope for either sign of the
centrifugal term and so cannot detect an error in the rotating frame. The
uniformly rotating sphere, whose slope field is given in closed form by
Eq.~\eqref{eq:rotsphere} and derived in Appendix~\ref{app:rotsphere},
discriminates sharply---the two signs differ by more than seven degrees at
mid-latitude for $\omegarot^{2}R/g = 0.5$. We recommend it as a routine check
and provide the implementation.

\end{enumerate}

Two qualifications frame these results. For bodies without a measured mass the
assumed bulk density remains the dominant uncertainty, shifting the mean slope
by some $10^{\circ}$ on Bennu across a plausible range---three times the effect
of mesh resolution---so that constraining the mass buys more than refining the
shape. And all fields assume a homogeneous interior, an assumption that leaves
the resolution results untouched but that is known to fail for bilobed bodies
such as Itokawa, whose head and body differ appreciably in density.

That last limitation points to the natural continuation. A surface shaped by
long-term downslope transport should approach an equipotential of the true
gravity field, so that the spread of geopotential over the surface becomes an
estimator of the interior density distribution rather than merely a description
of it. Such an inversion requires a diagnostic whose response to the assumed
density is not swamped by its response to the mesh---the property established
here for the geopotential dispersion and shown to be absent for the mean slope.
That development is reserved for future work and lies beyond the scope of the
present study.
% ============================================================
% DATA AVAILABILITY STATEMENT  --  SDGA Paper IV
%
% Paste this block at the END of Section8.tex, after the final
% paragraph of the Conclusions and before \end{document} is reached.
%
% The reserved DOI is already filled in below (10.5281/zenodo.21793575).
% No further editing is needed.
% ============================================================

\section*{Data availability}

The analysis code, the verification scripts, the five figures, and a
machine-readable file containing every numerical result quoted in this paper
are openly available at
\href{https://doi.org/10.5281/zenodo.21793575}{doi:10.5281/zenodo.21793575},
released under CC BY 4.0. The deposit includes the two verification programs
that implement Eqs.~\eqref{eq:ellipsoid} and~\eqref{eq:rotsphere}, so that both
checks can be reproduced independently of the shape models. Development of the
code is tracked at
\url{https://github.com/praponkan/shape-based-surface-gravity}.

Shape models are not redistributed. Those used here are available from the
OSIRIS-REx mission archive (Bennu), the JAXA DARTS archive accompanying
\citet{kanamaru2019} (Ryugu shape model and reference gravitational products),
and the PDS Small Bodies Node (Eros, Itokawa, Gaspra, Ida, Kleopatra); the
specific models, assumed densities, and rotation periods are listed in
Table~\ref{tab:shapes}.


% ============================================================
% OPTIONAL -- if the journal requires a separate acknowledgement
% of the software used, this can follow the statement above.
% ============================================================

% \section*{Software}
%
% Self-gravitation was computed with \texttt{polyhedral-gravity}
% version~3.3.1 \citep{gravitylib}. Analysis used Python~3.13 with
% \texttt{numpy} and \texttt{scipy}; figures were produced with
% \texttt{matplotlib}.




\appendix
\section{Gravitational slope on a uniformly rotating sphere}
\label{app:rotsphere}

This appendix derives the closed-form slope field used as a verification test in
Section~\ref{sec:numval}, and shows explicitly how it discriminates between the
two possible signs of the centrifugal term.

Consider a homogeneous sphere of radius $R$ and mass $M$, rotating uniformly at
rate $\omegarot$ about the $z$ axis, and write $g = GM/R^{2}$ for the magnitude
of the surface attraction. Because the configuration is axisymmetric, take
without loss of generality a surface point at longitude zero and latitude
$\phi$,

\begin{equation}
\mathbf{r} \;=\; R\,(\cos\phi,\;0,\;\sin\phi),
\qquad
\nhat \;=\; (\cos\phi,\;0,\;\sin\phi),
\label{eq:apppoint}
\end{equation}

the outward normal of a sphere being radial. The self-gravitational attraction
of a homogeneous sphere at its surface is radial and inward,

\begin{equation}
\ngrav \;=\; -\,g\,(\cos\phi,\;0,\;\sin\phi),
\label{eq:appattract}
\end{equation}

while the component of $\mathbf{r}$ perpendicular to the spin axis is
$\mathbf{r}_{\perp} = R\cos\phi\,(1,0,0)$, so that from Eq.~\eqref{eq:acf}

\begin{equation}
\acf \;=\; \omegarot^{2} R\cos\phi\,(1,\;0,\;0).
\label{eq:appcf}
\end{equation}

Introduce the dimensionless ratio of centrifugal to gravitational acceleration
at the equator,

\begin{equation}
k \;\equiv\; \frac{\omegarot^{2}R}{g},
\label{eq:appk}
\end{equation}

For every body considered in this paper $k$ is well below unity, and $k<1$ is
assumed in what follows unless stated otherwise; the case $k>1$ is treated
explicitly below. Adding Eqs.~\eqref{eq:appattract} and~\eqref{eq:appcf} gives
the effective acceleration

\begin{equation}
\geff \;=\; g\bigl(\,(k-1)\cos\phi,\;\;0,\;\;-\sin\phi\,\bigr).
\label{eq:appgeff}
\end{equation}

The slope of Eq.~\eqref{eq:slope} is the angle between $-\geff$ and $\nhat$.
Forming the inner product,

\begin{equation}
-\geff\cdot\nhat
\;=\; g\bigl[(1-k)\cos^{2}\phi + \sin^{2}\phi\bigr]
\;=\; g\bigl(1 - k\cos^{2}\phi\bigr),
\label{eq:appdot}
\end{equation}

and the magnitude,

\begin{equation}
\lVert\geff\rVert \;=\; g\sqrt{(1-k)^{2}\cos^{2}\phi + \sin^{2}\phi},
\label{eq:appmag}
\end{equation}

so that

\begin{equation}
\cos\slope(\phi) \;=\;
\frac{1-k\cos^{2}\phi}
     {\sqrt{(1-k)^{2}\cos^{2}\phi + \sin^{2}\phi}},
\label{eq:approtsphere}
\end{equation}

which is Eq.~\eqref{eq:rotsphere}. The factor $g$ cancels, so the slope field
depends on the body only through $k$.

Two limits check the result. At the pole, $\phi = 90^{\circ}$, both numerator
and denominator reduce to unity and $\slope = 0$: the spin axis passes through
the point, the centrifugal term vanishes there, and the effective gravity is
radial. At the equator, $\phi = 0$, the expression becomes $(1-k)/\lvert
1-k\rvert$, which is $+1$ for $k<1$ and again gives $\slope = 0$, the
centrifugal term being antiparallel to the attraction and so changing its
magnitude but not its direction. For $k>1$ the same expression gives $-1$ and
hence $\slope = 180^{\circ}$: the effective acceleration points outward and
loose material is no longer bound, which is the rotational disruption limit. The
slope is therefore zero at both extremes of latitude and attains its maximum in
between, which is what makes intermediate latitudes the informative ones for
verification.

The discriminating power of the test follows from the appearance of $k$ in
Eq.~\eqref{eq:approtsphere}. Reversing the sign of $\acf$ in
Eq.~\eqref{eq:geff}---adding the centrifugal acceleration to the attraction
rather than opposing it---is equivalent to replacing $k$ by $-k$ throughout, so
that the predicted slope becomes

\begin{equation}
\cos\slope_{\text{wrong}}(\phi) \;=\;
\frac{1+k\cos^{2}\phi}
     {\sqrt{(1+k)^{2}\cos^{2}\phi + \sin^{2}\phi}}.
\label{eq:appwrong}
\end{equation}

At $k = 0.5$ the two expressions give the values in
Table~\ref{tab:approtsphere}. They differ by $10.2^{\circ}$ at latitude
$30^{\circ}$ and by $7.1^{\circ}$ at $45^{\circ}$---differences far larger than
the discretisation error of a moderately refined mesh, so the test is decisive
rather than marginal. Both expressions vanish at $\phi = 0$ and
$\phi = 90^{\circ}$, which is precisely why a test confined to a non-rotating
body, or one that examined only the equator and poles, would detect nothing.

\begin{table}[t]
\centering
\small
\caption{Gravitational slope on a uniformly rotating sphere at
$k = \omegarot^{2}R/g = 0.5$, for the correct sign of the centrifugal term
(Eq.~\ref{eq:approtsphere}) and for the reversed sign
(Eq.~\ref{eq:appwrong}), together with the values returned by the
implementation on an icosphere mesh.}
\label{tab:approtsphere}
\begin{tabular}{rrrr}
\toprule
$\phi$ & Eq.~\eqref{eq:approtsphere} & Eq.~\eqref{eq:appwrong} & This work \\
\midrule
$30^{\circ}$ & $19.11^{\circ}$ & $8.95^{\circ}$  & $18.97^{\circ}$ \\
$45^{\circ}$ & $18.43^{\circ}$ & $11.31^{\circ}$ & $18.33^{\circ}$ \\
$60^{\circ}$ & $13.90^{\circ}$ & $10.89^{\circ}$ & $13.91^{\circ}$ \\
\bottomrule
\end{tabular}
\end{table}

The residual between the computed and analytic values, at most $0.14^{\circ}$,
is consistent with the polygonal approximation of the sphere at the mesh
resolution used and decreases under refinement.



\begin{thebibliography}{99}

% ---- companion papers (SDGA series) ----
\bibitem[Kanjanatarayont(2026a)]{paper1}
Kanjanatarayont, P., 2026. Surface gravity of small bodies from the
area-equivalent radius: an $e^{2}$-cancellation theorem and an
{$\eta \approx 1.05$} domain of validity. Companion paper, Paper I.
doi:10.5281/zenodo.21327611.

\bibitem[Kanjanatarayont(2026b)]{paper2}
Kanjanatarayont, P., 2026. The optimal area--volume radius blend for the
mean surface gravity of small bodies: a non-analytic $\Pthree/\Ptwo$ law.
Companion paper, Paper II. doi:10.5281/zenodo.21328178.

\bibitem[Kanjanatarayont(2026c)]{paper3}
Kanjanatarayont, P., 2026. Shape-based estimation of mean surface gravity
for small bodies: the accuracy of the blended radius and the limits of
convex models. Companion paper, Paper III. doi:10.5281/zenodo.21622962.

% ---- polyhedral gravity method ----
\bibitem[Werner \& Scheeres(1996)]{werner1996}
Werner, R.~A., Scheeres, D.~J., 1996. Exterior gravitation of a
polyhedron derived and compared with harmonic and mascon gravitation
representations of asteroid 4769 Castalia. \emph{Celestial Mechanics
and Dynamical Astronomy} 65, 313--344. doi:10.1007/BF00053511.

\bibitem[Tsoulis(2012)]{tsoulis2012}
Tsoulis, D., 2012. Analytical computation of the full gravity tensor
of a homogeneous arbitrarily shaped polyhedral source. \emph{Geophysics}
77, F1--F11. doi:10.1190/geo2010-0334.1.

\bibitem[Schuhmacher et~al.(2024)]{gravitylib}
Schuhmacher, J., Blazquez, E., Gratl, F., Izzo, D., G\'omez, P., 2024.
Efficient polyhedral gravity modeling in modern C++ and Python.
\emph{Journal of Open Source Software} 9(98), 6384.
doi:10.21105/joss.06384. Software: \texttt{polyhedral-gravity}, version~3.3.1.

% ---- surface slope / geopotential on small bodies ----
\bibitem[Scheeres et~al.(2016)]{scheeres2016}
Scheeres, D.~J., et al., 2016. The geophysical environment of Bennu.
\emph{Icarus} 276, 116--140. doi:10.1016/j.icarus.2016.04.013.

\bibitem[Scheeres et~al.(2019)]{scheeres2019}
Scheeres, D.~J., et al., 2019. The dynamic geophysical environment of
(101955) Bennu based on OSIRIS-REx measurements. \emph{Nature Astronomy}
3, 352--361. doi:10.1038/s41550-019-0721-3.

\bibitem[Barnouin et~al.(2019)]{barnouin2019}
Barnouin, O.~S., et al., 2019. Shape of (101955) Bennu indicative of a
rubble pile with internal stiffness. \emph{Nature Geoscience} 12,
247--252. doi:10.1038/s41561-019-0330-x.

\bibitem[Barnouin et~al.(2020)]{barnouin2020}
Barnouin, O.~S., et al., 2020. Digital terrain mapping by the
OSIRIS-REx mission. \emph{Planetary and Space Science} 180, 104764.
doi:10.1016/j.pss.2019.104764.

\bibitem[Daly et~al.(2020)]{daly2020}
Daly, M.~G., et al., 2020. Hemispherical differences in the shape and
topography of asteroid (101955) Bennu. \emph{Science Advances} 6,
eabd3649. doi:10.1126/sciadv.abd3649.

\bibitem[Jawin et~al.(2020)]{jawin2020}
Jawin, E.~R., et al., 2020. Global patterns of recent mass movement on
asteroid (101955) Bennu. \emph{Journal of Geophysical Research: Planets}
125, e2020JE006475. doi:10.1029/2020JE006475.

% ---- Eros ----
\bibitem[Zuber et~al.(2000)]{zuber2000}
Zuber, M.~T., et al., 2000. The shape of 433 Eros from the NEAR-Shoemaker
laser rangefinder. \emph{Science} 289, 2097--2101.
doi:10.1126/science.289.5487.2097.

\bibitem[Thomas et~al.(2002)]{thomas2002}
Thomas, P.~C., et al., 2002. Eros: shape, topography, and slope processes.
\emph{Icarus} 155, 18--37. doi:10.1006/icar.2001.6755.

\bibitem[Yeomans et~al.(2000)]{yeomans2000}
Yeomans, D.~K., et al., 2000. Radio science results during the
NEAR-Shoemaker spacecraft rendezvous with Eros. \emph{Science}
289, 2085--2088. doi:10.1126/science.289.5487.2085.

% ---- Itokawa + density estimation ----
\bibitem[Gaskell et~al.(2008)]{gaskell2008}
Gaskell, R.~W., et al., 2008. Characterizing and navigating small bodies
with imaging data. \emph{Meteoritics \& Planetary Science} 43, 1049--1061.
doi:10.1111/j.1945-5100.2008.tb00692.x.

\bibitem[Fujiwara et~al.(2006)]{fujiwara2006}
Fujiwara, A., et al., 2006. The rubble-pile asteroid Itokawa as observed
by Hayabusa. \emph{Science} 312, 1330--1334. doi:10.1126/science.1125841.

\bibitem[Abe et~al.(2006)]{abe2006}
Abe, S., et al., 2006. Mass and local topography measurements of Itokawa
by Hayabusa. \emph{Science} 312, 1344--1349. doi:10.1126/science.1126272.

\bibitem[Kanamaru \& Sasaki(2019)]{kanamaru2019}
Kanamaru, M., Sasaki, S., 2019. Estimation of interior density distribution
for small bodies: the case of asteroid Itokawa. \emph{Transactions of the
Japan Society for Aeronautical and Space Sciences, Aerospace Technology
Japan} 17, 270--275. doi:10.2322/tastj.17.270.

\bibitem[Lowry et~al.(2014)]{lowry2014}
Lowry, S.~C., et al., 2014. The internal structure of asteroid (25143)
Itokawa as revealed by detection of YORP spin-up. \emph{Astronomy \&
Astrophysics} 562, A48. doi:10.1051/0004-6361/201322549.

\bibitem[Ostro et~al.(2005)]{ostro2005}
Ostro, S.~J., et al., 2005. Radar observations of Itokawa in 2004 and
improved shape estimation. \emph{Meteoritics \& Planetary Science} 40,
1563--1574. doi:10.1111/j.1945-5100.2005.tb00136.x.

% ---- rubble piles / general ----
\bibitem[Walsh(2018)]{walsh2018}
Walsh, K.~J., 2018. Rubble pile asteroids. \emph{Annual Review of
Astronomy and Astrophysics} 56, 593--624.
doi:10.1146/annurev-astro-081817-052013.

\bibitem[Dobrovolskis(2019)]{dobrovolskis2019}
Dobrovolskis, A.~R., 2019. Classification of ellipsoids by shape
and surface gravity. \emph{Icarus} 321, 891--928.
doi:10.1016/j.icarus.2018.11.023.

\end{thebibliography}


\end{document}
